% =====================================================================
%  Capitolo 8 — Introduzione al level design: chi è il level designer,
%  quali sono i suoi strumenti, e i primi passi nell'editor di Unreal
%  Engine (primo anno)
%  Fonte: trascrizione della lezione GD8, che è in realtà la prima
%  lezione del corso di Level Design del primo anno (docente: Stefano
%  Zocchi), lezione ibrida (in aula a Milano e da remoto).
%  Contenuti, nell'ordine della lezione: annunci pratici; chi è il
%  docente; il giro di presentazioni (le sei domande, il livello
%  peggiore di ciascuno: la Cattedrale del Sangue Sacro di Code Vein,
%  Subnautica, Nightreign, la Città Infame di Dark Souls, Pentiment,
%  Slime Rancher, Frostpunk 2...); perché presentarsi: parlare in
%  pubblico; che cos'è il level design; che cosa fa un level designer
%  (blockout, luci, set dressing, documentazione, test, narrazione
%  ambientale); qual è lo strumento del level designer ("dipende", e la
%  risposta vera: sé stessi, imparare a imparare); teoria, pratica,
%  expertise, esperienza sul campo, lavoro di gruppo; il picco della
%  fiducia e l'effetto Dunning-Kruger; i tre corollari (niente regole
%  fisse, design olistico, design come artigianato); la struttura del
%  corso; la pausa; il progetto in Unreal Engine 5; l'editor
%  (viewport, Place Actors, gizmo, Outliner, Content Browser, asset e
%  actor, transform, Details, scorciatoie W/E/R, gruppi, istanze);
%  la documentazione; i livelli e il salvataggio; le scorciatoie
%  utili (snap to floor, operazioni nei campi numerici); il task: il
%  blockout della propria stanza; il materiale extra (GDC, An
%  Architectural Approach to Level Design); la chiusura.
%  Integrazioni dal libro "An Architectural Approach to Level Design"
%  (Christopher W. Totten), segnalate come approfondimenti.
%
%  Pacchetti richiesti nel preambolo del libro (già presenti nel
%  main.tex): \usepackage{booktabs}, \usepackage{amsmath},
%  \usepackage{tikz} con
%    \usetikzlibrary{shapes.geometric,shapes.symbols,shapes.misc,
%      arrows.meta,positioning,calc,fit}
%  (le librerie vanno caricate nel preambolo, non nel capitolo).
%  Tutte le etichette di questo capitolo hanno il prefisso "gd8:"
%  (gd8:sec:..., gd8:fig:..., gd8:tab:...) e tutti gli stili TikZ il
%  prefisso "gd8", per non collidere con gli altri capitoli.
%
%  Nota sulla fonte: la trascrizione automatica è, come per GD7,
%  molto rovinata. L'inizio della registrazione, la pausa e i tratti di
%  lavoro individuale sono stati trascritti come frasi senza senso o
%  ripetute in loop ("non funziona", "mi fai il mio sguardo", "un altro
%  mixer"...), con interi brani captati da una fonte esterna (una radio
%  o un video: frasi sul "mercato del lavoro", sulla "Corte Suprema di
%  Internet", sulla "struttura del giardino", sugli "egiziani"), che qui
%  sono omessi. Durante la pausa la registrazione cattura anche
%  conversazioni private degli studenti (su un gioco di carte e sul suo
%  meta), che non fanno parte della lezione e non sono riportate. Molti
%  nomi propri (di giochi, persone, strumenti) sono storpiati: dove il
%  riferimento è ricostruibile con certezza lo si riporta nella forma
%  corretta; dove non lo è, lo si segnala in nota invece di indovinare.
%  Le figure sono ricostruzioni schematiche di ciò che è mostrato a
%  lezione, non le slide originali.
% =====================================================================

\chapter[Introduzione al level design]{Introduzione al level design: lo
strumento sei tu, l'engine viene dopo}
\label{gd8:cap:leveldesign}

Il capitolo precedente ha aperto Unreal Engine dal lato del
\textbf{codice}: eventi, nodi, variabili (capitolo~\ref{gd7:cap:blueprint}).
Questo capitolo lo riapre dal lato dello \textbf{spazio}. La lezione da cui
nasce è infatti la \textbf{prima lezione del corso di Level Design} del primo
anno, tenuta da un docente nuovo, \textbf{Stefano Zocchi}, game e level
designer: ``benvenuti alla prima lezione, per il primo anno, nel corso di Game
Design, specificatamente alla lezione del corso di \textbf{Level Design}''.

Il corso di level design, nel percorso di Game Design, è uno dei tre filoni
che corrono in parallelo (gli studenti stessi parlano delle ``tre materie che
dobbiamo fare nel corso''), e ha un taglio preciso: \textbf{molta pratica in
engine}, accompagnata da teoria e metodologia. Questa prima lezione, però, è
soprattutto un'apertura: chi è il docente, chi sono gli studenti, che cos'è il
level design, che cosa fa davvero un level designer, e qual è l'unico
strumento di cui ha veramente bisogno. Solo nell'ultima ora si apre l'editor,
per imparare a muoversi, piazzare forme e organizzare una scena; e il primo
task è il più semplice e personale possibile: ricostruire con le forme di base
\textbf{la propria stanza}.

Chi ha seguito il corso breve su Unreal troverà la parte pratica in gran parte
nota (capitolo~\ref{ue1:cap:intro}, in particolare le
sezioni~\ref{ue1:sec:interfaccia} e~\ref{ue1:sec:transform}); qui la si
ripercorre con l'occhio del level designer, che passa ``quasi tutto il tempo
in engine a piazzare cose, sistemare cose, testare cose''. Il blockout vero e
proprio, con le luci, verrà poi ripreso in modo approfondito nel corso su
Unreal, nel capitolo dedicato alla \emph{gym}, al blockout e alle luci.

La struttura della lezione, e quindi del capitolo, è questa:
\begin{enumerate}
  \item gli \textbf{annunci pratici} e la presentazione del docente
        (sezioni~\ref{gd8:sec:annunci} e~\ref{gd8:sec:docente});
  \item il \textbf{giro di presentazioni}: sei domande a ciascuno, tra cui
        il livello peggiore che si sia mai giocato e come lo si
        cambierebbe (sezione~\ref{gd8:sec:presentazioni});
  \item perché presentarsi: la prima lezione seria è \textbf{saper parlare
        in pubblico} (sezione~\ref{gd8:sec:pubblico});
  \item \textbf{che cos'è} il level design e \textbf{che cosa fa} un level
        designer (sezioni~\ref{gd8:sec:cos-e} e~\ref{gd8:sec:cosa-fa});
  \item lo \textbf{strumento} del level designer, la parola magica
        ``dipende'' e l'imparare a imparare
        (sezione~\ref{gd8:sec:strumento});
  \item come si diventa level designer: teoria, pratica, expertise,
        esperienza, gruppo (sezione~\ref{gd8:sec:imparare});
  \item due curve da conoscere in anticipo: il \textbf{picco} della
        fiducia e l'\textbf{effetto Dunning-Kruger}
        (sezione~\ref{gd8:sec:curve});
  \item i \textbf{tre corollari}: niente regole fisse, il design è olistico,
        il design è artigianato (sezione~\ref{gd8:sec:corollari});
  \item la \textbf{struttura del corso} e la pausa
        (sezione~\ref{gd8:sec:struttura});
  \item \textbf{Unreal Engine}: il progetto, l'editor, gli actor e gli asset,
        il transform, le scorciatoie, i livelli
        (sezioni~\ref{gd8:sec:progetto}--\ref{gd8:sec:livelli});
  \item il \textbf{task}: il blockout della propria stanza, e il materiale
        extra per chi vuole approfondire
        (sezioni~\ref{gd8:sec:task} e~\ref{gd8:sec:extra}).
\end{enumerate}

% ---------------------------------------------------------------------
\section{Prima di cominciare: gli annunci pratici}
\label{gd8:sec:annunci}

\subsection{Cinque minuti di margine}

Il docente apre con una ``piccola informazione di servizio'': lascia
\textbf{sempre cinque minuti} prima di iniziare, per dare a tutti il tempo di
arrivare. Capita spesso, spiega, che ci sia ``gente che viene da fuori, o che
c'è gente che lavora''; e se ci sono autobus in ritardo o ``casini vari'', è
giusto che chi arriva possa sedersi con calma. Solo dopo si fanno
presentazioni, introduzioni e tutto il resto.

\subsection{In aula o da remoto}

Le lezioni, ``come sapete perfettamente'', si possono seguire
\textbf{contemporaneamente} in aula e online. Il docente non nasconde la sua
preferenza: in classe è ``l'ambiente migliore'', perché è più veloce parlare
con le persone, far vedere le cose, \textbf{andare ai vostri computer} e
verificare se c'è qualcosa da sistemare. Ma se serve seguire da remoto, ``per
me non c'è nessun problema''. L'unica richiesta: se c'è un motivo particolare,
avvisarlo prima se non si è in grado di venire. Per il resto, ``la praticità
conta più di tutto''.

\subsection{Pausa e appello}

Il ritmo della lezione è sempre lo stesso:
\begin{itemize}
  \item qualche minuto di margine all'inizio;
  \item a \textbf{metà lezione} una \textbf{pausa} di 15, 20 o 30 minuti,
        ``a seconda di quanto è stato pesante l'argomento'' trattato fino a
        quel momento;
  \item \textbf{dopo la pausa}, l'\textbf{appello}. Non lo fa subito
        proprio per segnare come presente chi arriva in ritardo. Chi arriva
        dopo l'appello deve farglielo sapere, così lo segna lo stesso; e se
        il docente si dimentica di farlo, va ricordato: ``mi è capitato più
        di una volta. Siamo umani''.
\end{itemize}

\paragraph{Un gruppo nuovo.} Il primo anno di quest'anno è ``più numeroso di
quello dell'anno scorso, fantastico''; e, per la prima volta nella sede di
Milano, ci sono anche delle ragazze a lezione, ``un'ottima cosa''.

% ---------------------------------------------------------------------
\section{Chi insegna}
\label{gd8:sec:docente}

\paragraph{Regola numero uno.} ``Regola numero uno dell'apprendimento:
\textbf{sapere sempre chi vi sta insegnando qualcosa}''. Mai fidarsi della
gente a caso su internet, ``soprattutto su YouTube'' (un tema che il docente
promette di riprendere durante l'anno).

\paragraph{La carriera.} Stefano Zocchi è un \textbf{game e level designer}.
Il suo percorso, in breve:
\begin{itemize}
  \item è partito dal mondo \textbf{indie}, partecipando a vari progetti,
        facendo ``un po' di roba horror per i fatti miei'' e lavorando con
        altre persone;
  \item è stato poi assunto da uno studio gallese, \textbf{Antimatter
        Games}, dove ha lavorato a \textbf{sparatutto stealth strategici}.
        Lo studio è stato chiuso durante l'ondata di chiusure che ha colpito
        l'industria dei videogiochi negli ultimi anni;
  \item ha lavorato per un periodo per \textbf{Behaviour Interactive}, nome
        noto soprattutto per \emph{Dead by Daylight} (il gioco horror
        multiplayer in cui un giocatore interpreta il killer e gli altri
        quattro i sopravvissuti; il titolo glielo suggerisce la classe).
        Behaviour, però, è ``un ombrello di un sacco di studi diversi'', e
        lui non lavorava a quel gioco. A che cosa lavorasse, non può dirlo;
  \item oggi lavora come \textbf{freelance}: insegna, collabora con un paio
        di ex studenti e segue vari progetti.
\end{itemize}

\paragraph{L'NDA.} Il motivo per cui non può dire a che cosa ha lavorato è
un \textbf{NDA}, un \emph{non-disclosure agreement}: una parte del contratto
che si firma e che stabilisce che non si può parlare di ciò che lo studio non
autorizza esplicitamente. C'è una ``grossa parte'' del suo curriculum per cui
può dire ``ho lavorato in questo studio'', ma alla domanda ``bello, cosa hai
fatto?'' la risposta è sempre la stessa: ``vorrei continuare a lavorare
nell'industria dei videogiochi''. È una realtà del mestiere con cui chiunque,
nel corso della carriera, si troverà a fare i conti.

\paragraph{Il progetto attuale.} Il progetto più grosso a cui sta lavorando
si chiama \textbf{\emph{Break the Night}}: un gioco d'azione corpo a corpo a
tema un po' \emph{Y2K}, un po' \emph{pop punk}, ``energico'', con ``un sacco
di combo'', attacchi pesanti, attacchi leggeri, mosse speciali, cancellazioni
con gli scatti.\footnote{La trascrizione definisce il genere con una parola
storpiata (``una cashlash''); dalla descrizione (combo, attacchi leggeri e
pesanti, mosse speciali, \emph{cancel} con lo scatto) si tratta di un
\emph{action} corpo a corpo, verosimilmente un \emph{hack and slash} o un
\emph{beat 'em up}, ma il termine esatto usato dal docente non è
ricostruibile.} ``Molto, molto divertente.''

\paragraph{Come è arrivato a insegnare.} Qualche anno fa un insegnante non
poteva più seguire il corso di Milano; l'altro docente che si occupa di questi
corsi gli ha chiesto se poteva coprirlo lui. Essendo freelance, si inseriva
``abbastanza bene'' nei lavori che stava facendo: bastava riorganizzare un po'
gli orari. Da qualche anno, ormai, insegna lui il corso di level design.

\paragraph{Alessandro.} In aula c'è anche \textbf{Alessandro}, che non fa
parte del primo anno, ma che gli studenti vedranno ``un po' dappertutto in
questo edificio, come un fantasma''.\footnote{La trascrizione aggiunge un
dettaglio incomprensibile (``che va zika per via macchina''); non lo si
ricostruisce.}

% ---------------------------------------------------------------------
\section{Il giro di presentazioni}
\label{gd8:sec:presentazioni}

``Basta con queste frasi noiose'': è un anno nuovo, il docente non conosce
ancora la classe, e vuole sentire tutti, in aula e online. Le domande vanno
``dal semplice al più interessante'', per capire che tipo di persona è
ciascuno.

\subsection{Le sei domande}

\begin{enumerate}
  \item \textbf{Chi siete?}
  \item \textbf{Che tipo di giochi volete fare?}
  \item \textbf{Qual è la vostra forza principale?}
  \item \textbf{Qual è la vostra debolezza principale?}
  \item \textbf{In che cosa vi vorreste specializzare}, se già avete
        un'idea?
  \item E, visto il tema del corso: \textbf{un livello che vi ha proprio
        fatto schifo}, che vi ha fatto mollare il gioco, ``che ancora ce
        l'avete qua'' e vi chiedete come sia possibile che abbiano fatto
        ``una roba così orrenda''. \textbf{Perché} non vi è piaciuto, e
        \textbf{che cosa fareste per migliorarlo}?
\end{enumerate}

L'ultima domanda è quella che conta per il corso: non basta dire che un
livello è brutto, bisogna saper dire \emph{perché} e proporre un
\emph{rimedio}. È, in piccolo, il lavoro di analisi che il level designer fa
ogni giorno.

\subsection{L'esempio del docente: la Cattedrale del Sangue Sacro}

Siccome la prima persona a rispondere ``si trova un attimo sulle spine'',
il primo a rispondere all'ultima domanda è il docente stesso. Ogni anno rifà
le slide; questa, per la prima volta, non l'ha rifatta, perché il livello in
questione ce l'ha in testa da quando ci ha giocato due anni fa e ``continua a
tormentarmi i sogni'': gli appare ogni volta che sta facendo un livello,
dicendogli ``vuoi fare questa roba? Non ti libererai mai di me''.

Il livello è la \textbf{Cattedrale del Sangue Sacro} di \textbf{\emph{Code
Vein}}, uno dei ``tantissimi giochi'' usciti sulla scia dei titoli di
FromSoftware, un \emph{souls-like} ``tra virgolette''.\footnote{La
trascrizione riporta prima ``sulla scia di Vlaovic'', poi si corregge in
``sulla scia di From Software''; il primo nome non è ricostruibile.} Per il
docente è uno degli esempi che toccano ``il punto più fondo di tutto questo
abisso di mediocrità''.

La slide mostra la mappa del livello, realizzata da un utente di DeviantArt
per chi non riusciva a capire dove andare.\footnote{Il nome dell'utente
ringraziato dal docente è storpiato nella trascrizione (``Cercheron
Quartoro'') e non è ricostruibile.} Le linee grigie non sono strade: sono
\textbf{passerelle}. La domanda alla classe è: guardando questa mappa, in
un'ottica negativa, che cosa salta all'occhio? La risposta arriva subito:
\textbf{è un labirinto}. Il docente elenca i problemi:
\begin{description}
  \item[Un labirinto di corridoi.] È ``una lunghissima serie di
        corridoi'', con bivi di cui non si capisce la direzione giusta;
        corridoi che ``vanno in giro totalmente a caso''.
  \item[Su più livelli, non collegati.] I corridoi non sono tutti
        collegati fra loro: stanno su piani diversi. Se ci si trova in un
        punto e si vuole raggiungerne un altro, magari bisogna tornare
        indietro, scendere, prendere un altro corridoio e girare a destra;
        ``e io non lo so''.
  \item[Un ``pastone'' visivo.] Tutti i corridoi sono circondati da enormi
        decorazioni gotiche, ``un po' Duomo di Milano, un po' cattedrali
        inglesi'': archi, pilastri, ``rampart'', decorazioni a caso. Rendono
        impossibile capire anche solo dove si trova il corridoio successivo.
  \item[Nessun punto di riferimento.] Non si capisce dove sono i punti di
        riferimento, non si capisce dove andare anche quando si intuisce più
        o meno la direzione.
  \item[Ripetitivo e geometrico.] È ``ripetitivo da far paura'' e, come
        mostra la struttura della mappa, totalmente geometrico: dall'inizio
        alla fine è sempre lo stesso schema, corridoio e torre, corridoio e
        torre. \textbf{Non c'è interesse nell'esplorazione}.
  \item[E i boss sono terribili.] Risultato: ``ho mollato il gioco'' in quel
        punto.
\end{description}
È, in negativo, una lista di quasi tutti i temi del corso: navigazione,
leggibilità, punti di riferimento (\emph{landmark}), varietà, interesse
dell'esplorazione.

\subsection{Valerio: Subnautica}

Il primo studente è \textbf{Valerio}. Non sa rispondere a due domande: che
tipo di giochi vuole fare e in che cosa vorrebbe specializzarsi. ``Ci sta. È
abbastanza comune questa risposta al primo anno'': si sta iniziando adesso, e
probabilmente l'unica cosa che accomuna tutti i presenti è che gli piacciono i
videogiochi (``c'è qualcuno a cui non piacciono i videogiochi?'').

Più che un livello, Valerio cita un gioco intero che gli ha fatto smettere di
giocare: \textbf{\emph{Subnautica}}, perché ``non aveva troppa libertà'' e non
sapeva come muoversi; e, in generale, non ama i giochi di sopravvivenza e di
esplorazione. Gliel'aveva consigliato un amico appassionato. Come debolezza,
dice di fare fatica a impegnarsi in una cosa se non lo prende
particolarmente; come forza, non sa. ``Molto onesto, molto diretto. Apprezzo
questo approccio.'' Alla domanda su che cosa cambierebbe di \emph{Subnautica},
la risposta è la più breve della giornata: ``cancellerei il gioco''. Il
docente apprezza il tempismo comico.

\paragraph{I nomi.} Il docente avverte che ripeterà spesso i nomi di tutti,
perché fa ``una fatica enorme'' a ricordarli, e si scusa in anticipo. E
aggiunge che, come tutte le esperienze educative, anche questa serve
``anche per scoprire chi siamo''.

\subsection{Alberto: Nightreign, e la differenza tra scelta ed errore}

\textbf{Alberto} vorrebbe fare ``un po' di tutto'': non ha un genere
preferito, dai giochi di calcio ai \emph{gacha} giapponesi.\footnote{La
trascrizione è qui in loop (``un po' di tutto'' ripetuto decine di volte) e
contiene uno scambio scherzoso su un simulatore di autobus nella Berlino degli
anni Quaranta, che non è chiaro a chi vada attribuito.} La sua forza è
l'\textbf{attenzione ai dettagli}: nota subito se in una scena c'è qualcosa
leggermente spostato, ``un pallino''. È un punto di forza, ``ma anche no.
Dipende dalla situazione''. La sua debolezza è il \textbf{multitasking}: non
riesce a concentrarsi su più cose nello stesso momento.

\paragraph{Il multitasking non esiste.} Il docente lo rassicura: non è
necessariamente una debolezza, perché, a livello di capacità biologica di
elaborare le informazioni, \textbf{il multitasking non esiste}. Quello che
esiste è spostare completamente l'attenzione da una cosa all'altra a velocità
straordinaria. Se ne riparlerà più avanti nel corso (il docente accenna a
``sistema 1 e sistema 2'').

\paragraph{Level o game?} Alberto è indeciso se specializzarsi in level design
o in game design, ``senza aver mai fatto una lezione di level design''. ``Ti
farò cambiare idea'', promette il docente.

\paragraph{Il livello: la mappa di Nightreign.} Il livello che non gli è
piaciuto viene da \textbf{\emph{Elden Ring Nightreign}}: una mappa (di un
contenuto aggiuntivo) con un castello da una parte e un altro dall'altra, in
cui bisogna risolvere delle cose su entrambi i lati per poi andare in un
altro castello ancora, al centro, e battere vari boss in successione. Ma non
c'è il tempo: come in un \emph{battle royale}, \textbf{la zona si stringe}, e
non si riescono a fare entrambi i castelli prima di arrivare a quello
centrale.

\paragraph{Scelta creativa o errore?} La domanda del docente è
provocatoria: ``secondo te, a FromSoftware sono così stupidi da non essersi
resi conto che un giocatore non può fare entrambi i castelli?'' No: se ne sono
resi conto, ma la scelta è voluta, perché il gioco è pensato per essere
giocato in gruppo, mandando qualcuno da una parte e qualcuno dall'altra. Qui,
dice il docente, ``iniziamo ad affrontare un discorso molto delicato e
complesso'': \textbf{quando una cosa che non ci piace è una scelta creativa, e
quando è un errore?} In questo caso diventa sempre più chiaro che impedire al
giocatore di fare entrambe le parti della mappa (o costringerlo a dividersi e
riunirsi al centro) è fatto apposta.

\paragraph{La proposta, e le sue conseguenze.} Alberto cambierebbe la cosa
lasciando la libertà di andare in quella zona anche quando si chiude: rendere
la parte fuori zona \textbf{non dannosa} per un certo tempo, così da poterla
completare anche dopo. Il docente fa notare che la modifica ``ha tutta una
serie di conseguenze'': prima di tutto è \textbf{antitetica all'idea di battle
royale}, in cui il campo di gioco si restringe proprio per costringere i
giocatori a scontrarsi. Alberto ribatte che lì lo scopo è esplorare e
potenziarsi il più possibile per battere il boss finale.

\paragraph{L'intervento di Mattia.} Interviene \textbf{Mattia}: più che un
vero battle royale (in cui, come dice il nome, alla fine ne resta uno solo),
\emph{Nightreign} è una \textbf{boss rush a tempo}, in cui bisogna arrivare a
un certo punto della mappa a seconda di come si restringe. Secondo lui il
fattore di level design è proprio questo: a meno di dividersi, o se si gioca
da soli, non si riesce a fare tutto, e questo serve a \textbf{impedire al
giocatore singolo di diventare troppo forte}. Del resto, aggiunge, non è un
gioco ``troppo divertente'' da soli. Il docente chiude la parentesi: la classe
è numerosa e c'è meno tempo per ``le mie teatralità''.

\subsection{Mattia: Monster Hunter, il sound design e la Città Infame}

\textbf{Mattia Maestri} non sa ancora che tipo di gioco vorrebbe creare, ma ha
un'idea dei giochi a cui gli piacerebbe lavorare: i \textbf{souls-like}, i
suoi generi preferiti, e, sogno ``nel cassetto'', entrare in
\textbf{Capcom} per lavorare a \textbf{\emph{Monster Hunter}}, la prima saga
che abbia mai giocato e la sua preferita (``gli ho giocati praticamente quasi
tutti''). Il docente nota la ``linea di discendenza'' interessante: i giochi di
FromSoftware prendono molto da \emph{Monster Hunter}.

\paragraph{Forza e debolezza.} La sua forza è la \textbf{creatività}: le idee
gli vengono facilmente (``che poi siano buone idee è un altro punto di
vista''). La debolezza: farà fatica nel level design, soprattutto nella parte
\textbf{tecnica}; si definisce ``più una persona teorica che una persona
pratica''.

\paragraph{Il sound design.} Oltre al game design, vorrebbe specializzarsi in
\textbf{sound design}: ama la musica, ``vivo di musica'', e alcuni giochi li
ha iniziati solo per la colonna sonora. \emph{Kingdom Hearts}, per esempio,
lo ha cominciato perché un amico gli aveva fatto sentire la musica del menu
iniziale.\footnote{Lo scambio sul brano è confuso: la classe cita un verso
(``If you walk away\dots''), che appartiene a una canzone di \emph{Kingdom
Hearts II}, e il docente lo identifica con un titolo che la trascrizione
riporta come ``Everly Hearts'', verosimilmente \emph{Dearly Beloved}, il tema
del menu della serie. Non è possibile stabilire con certezza a quale dei due
brani si riferisse Mattia.} All'inizio dell'anno aveva già chiesto se il sound
design sarebbe stato trattato, e gli era stato detto di approfondirlo fuori.

Il docente lo conferma: ci saranno ``delle piccole cose'' di sound design nel
corso, ``ma sono veramente poche''. Il consiglio migliore è un altro: nei
\textbf{progetti di gruppo} dei tre anni, quando il team si chiede ``che cosa
facciamo con quei suoni?'', bisogna essere quelli che dicono ``\textbf{ragazzi,
ci penso io}''. È così che ci si costruisce una specializzazione dentro un
percorso generalista.

\paragraph{Il livello: la Città Infame di Dark Souls.} Mattia ha platinato il
primo \emph{Dark Souls}, ma un livello gli è rimasto sullo stomaco: la
\textbf{Città Infame} (\emph{Blighttown}). È stata la parte che gli ha
richiesto più tempo, per tre motivi:
\begin{itemize}
  \item la \textbf{scarsa visibilità}, in una città tutta di legno e ``tutta
        uguale'';
  \item i \textbf{checkpoint}: solo due falò in tutta l'area, uno
        all'inizio e uno quando si è ``già praticamente fuori'', invece di
        tre o quattro intermedi;
  \item il \textbf{veleno}: quasi tutti i nemici lo infliggono, e si ha un
        tempo limitato per raggiungere il falò successivo, che però non si
        vede perché è lontanissimo.
\end{itemize}
Che cosa cambierebbe? Più visibilità, per capire meglio l'ambiente, oppure un
\textbf{falò intermedio}.

\paragraph{Fare souls-like.} Sui souls-like il docente è franco: da quando
sono usciti i giochi di FromSoftware, ``una quantità gigantesca'' di studenti
arriva dicendo ``io vorrei fare i souls-like'', e la risposta è sempre la
stessa: \textbf{pochissima gente riesce a farli}. Mattia precisa che intende
più in generale il genere \textbf{action RPG}. Non è impossibile: alcuni
studenti hanno portato come progetto finale un piccolo souls-like ``venuto
sorprendentemente bene''\footnote{Il titolo del progetto è storpiato nella
trascrizione (``Steam qualcosa'', ``Steam Rush'') e non è ricostruibile con
certezza.}; ovviamente non si può fare un gioco da sessanta ore nel contesto di
uno studio, ma il risultato era buono. E in Capcom qualche italiano c'è, anche
se quelli che il docente vede lavorare all'estero sono più nel dipartimento
artistico che in quello di design.

\subsection{Francesco: i JRPG, Fabula Ultima e il salto finale}

\textbf{Francesco Muratore} ama i \textbf{JRPG}: più \emph{Persona} che i
vecchi \emph{Final Fantasy}, anche se ha giocato di recente al settimo
capitolo; adesso sta giocando una saga ``non famosissima'', quella di
\emph{Trails}.\footnote{La trascrizione riporta ``trains in the sky''; dal
seguito del discorso (una serie con moltissimi capitoli collegati tra loro) si
tratta quasi certamente della serie \emph{The Legend of Heroes: Trails}, di
cui \emph{Trails in the Sky} è il primo capitolo.} Il docente conosce la
sensazione: a un certo punto ci si rende conto che una serie ``ha quindici
giochi'', tutti collegati. ``Giocare a certi JRPG è come giocare una
telenovela lunga vent'anni.''

Vorrebbe fare JRPG perché sono giochi con una storia che ``sa riempire'':
ogni volta che ne finisce uno resta appagato per un mese, un mese e mezzo,
prima di iniziarne un altro. Il docente traduce: ``un livello di passione
equivalente a quando ti rompi un ginocchio e non puoi andare a correre per un
mese e mezzo''.

\paragraph{Forza e debolezza.} La forza: ``mi rialzo facilmente dai miei
problemi e li affronto'', a testa alta. La debolezza: si distrae spesso, ma ci
sta lavorando. Il docente osserva che non è solo una questione di disciplina:
può essere una questione di \textbf{neurodivergenza}, di disturbi
dell'attenzione (``forse anche io''), e non c'è ``assolutamente niente di
male'': ``siamo tutti a gestirci il corpo che abbiamo, e il cervello è un
organo''.

\paragraph{Fare JRPG in Italia.} La reazione istintiva sarebbe ``buona
fortuna, fare JRPG in Italia''. Eppure c'è \textbf{\emph{Fabula Ultima}}: un
gioco di ruolo carta e penna italiano nato da chi era stanco del fantasy
classico europeo e americano di \emph{D\&D} e ha fatto un gioco da tavolo
basato sui JRPG, ``molto Final Fantasy'', come si capisce dal nome. Ha avuto
un successo ``stratosferico'': secondo il docente è probabilmente il gioco di
ruolo carta e penna italiano di maggior successo della storia.

\paragraph{Un paese di nippofili.} Del resto, ``viviamo in un paese di
giapponofili'': l'Italia è stata uno dei primi paesi al mondo a importare
anime, già dagli anni Ottanta, ben prima degli Stati Uniti. Con la
liberalizzazione dei canali, la televisione privata italiana aveva bisogno di
riempire il palinsesto, e i giapponesi vendevano serie da doppiare a prezzi
bassissimi: \emph{Heidi}, \emph{Sampei}, e poi \emph{I Cavalieri dello
Zodiaco}, \emph{Dragon Ball}.\footnote{La trascrizione cita anche un titolo
storpiato (``l'Upen''), forse \emph{Lupin III}; non è ricostruibile con
certezza.} Non è impossibile, quindi; è difficile, perché il mercato
occidentale è ``un po' recalcitrante'', ma ``viviamo in un mondo
globalizzato''. Il consiglio: \textbf{studiare i sistemi}, perché i JRPG sono
``tutti sistemi''. Il loro level design, sinceramente, ``tende a non essere
granché'': sono molto più concentrati su narrativa e sistemistica.

\paragraph{Il livello: un salto al 99,6\%.} Francesco cita un livello battuto
da poco in cui, al \textbf{99,6\%} del percorso, c'è un \textbf{salto finale
precisissimo}: ci è morto almeno sette volte, dopo cinque minuti di livello
tutti di fila.\footnote{Prima di questo esempio Francesco nomina un altro
titolo, irriconoscibile nella trascrizione (``su Giove ti lascia''); anche il
gioco del salto finale non è nominato in modo comprensibile.} È la classica
\textbf{botta di difficoltà} finale che costringe a rifare tutto da capo.
Taglierebbe quel pezzo? Il docente aggiunge una considerazione: forse proprio
quella è la sfida vera, la parte divertente. Anche questo va considerato.

\paragraph{L'inferno voluto.} Lo stesso vale per la Città Infame di Mattia:
è \textbf{pensata per essere l'inferno}, il punto più basso di tutto il mondo
di gioco. Una volta finita, la soddisfazione è ``non ci devo più tornare,
finalmente ce l'ho fatta''. È ``leggendaria'', uno dei livelli più difficili e
più ``tossici'' (``anche letteralmente'') di tutti gli action RPG. E c'è un
motivo se, scegliendo all'inizio la classe del ladro, si ha la
\textbf{Chiave Maestra}: permette di saltare l'intero livello, ed è ciò che
molti fanno alla seconda partita.

\subsection{Federica: i giochi intimi e la mappa di Pentiment}

\textbf{Federica} parte dalla debolezza, perché si collega al resto: una
grandissima \textbf{sensibilità}, ``in positivo, ma soprattutto in negativo'',
che la porta a rimuginare molto. Ed è anche ciò che vorrebbe fare: un
\textbf{piccolo gioco intimo}, come \emph{To the Moon}, \emph{Undertale},
\emph{Omori}. Il docente le consiglia \textbf{\emph{OFF}}: se le è piaciuto
\emph{Undertale}, è uno dei giochi a cui Toby Fox si è palesemente ispirato.

La sua forza è legata alle sue neurodivergenze, che a volte la aiutano a
\textbf{guardare le cose fuori dagli schemi} e la rendono molto creativa. Non
ha un'idea precisa di specializzazione: il suo desiderio è sempre stato
lavorare per conto suo, fare il suo piccolo gioco, anche senza fare molti
soldi, perché un lavoro ce l'ha già.

\paragraph{Il livello: la mappa di Pentiment.} Cita \textbf{\emph{Pentiment}},
che le stava piacendo ma ha lasciato. Il problema è la \textbf{mappa}, e
l'essere in \textbf{2D}: capisce la scelta di riprendere le miniature e il
disegno dell'epoca, ma quando si apre la mappa e ci si vuole spostare da un
punto all'altro si può andare ``a destra o a sinistra, a grandi linee''. Se si
deve raggiungere un punto molto lontano, lo spostamento è dispendioso; e la
mappa è poco chiara su quali NPC si trovino dove. La soluzione: un
\textbf{viaggio rapido}, o \textbf{più indicatori} sulla mappa. Più che un
livello specifico, nota il docente, è un problema di approccio alla
comunicazione. Come livello vero e proprio, Federica cita anche il
\textbf{Picco di Cristallo} di \emph{Hollow Knight}, che però mette alla prova
soprattutto la sua difficoltà con i \emph{platform}.

``Pochi masochisti hardcore gamer quest'anno'', commenta il docente
(l'anno prima c'era un appassionato di \emph{bullet hell}): Federica è una
delle persone ``con gli interessi più indie'' che abbia avuto finora.

\subsection{Daniele: l'esplorazione e gli spazi vuoti}

Si torna in aula, in senso antiorario, ``per tenere la gente sulle spine''.
\textbf{Daniele} vorrebbe fare giochi con un po' di
\textbf{esplorazione}, non necessariamente un \emph{open world}: anche
qualcosa di più piccolo, come esplorare una città.

\paragraph{Forza e debolezza.} La forza: è \textbf{preciso}, attento ai
dettagli ``senza fare le maniacate''. Il docente aggiunge una nota di
produzione: \textbf{c'è sempre una quantità di tempo limitata}, e l'attenzione
ai dettagli può diventare un problema quando si parla di tempistiche. ``Non si
può sempre fare tutto perfetto.'' La debolezza: non si sente tagliato per fare
il \textbf{leader} di un gruppo, preferisce non assumersi responsabilità. Il
docente: è importante saperlo. ``È bello dire: ah, io sono il capo''; ma
``vi assicuro di persona che essere il capo di un progetto non è divertente
per niente'', e non è sempre una buona idea essere quello che decide tutto e si
occupa di tutto.

\paragraph{Il livello: una montagna vuota.} Daniele descrive una zona di
montagna, in un gioco che gli è piaciuto, che però è un \textbf{grande vuoto}:
c'è poca roba, ci si arriva, ci si trova una grotta e poco
altro.\footnote{Il passaggio è quasi del tutto incomprensibile nella
trascrizione (``la parte della serena, papa, che è la montagna'') e il gioco
non è nominato in modo riconoscibile. Il docente collega l'esempio agli
``ultimi giochi di Pokémon'', e Daniele lo confronta con un'altra zona della
stessa mappa, un deserto: è quindi probabile, ma non certo, che si tratti di un
capitolo recente di \emph{Pokémon}.} Metterebbe \textbf{punti di interesse}:
come in \emph{Breath of the Wild}, dove su una montagna c'è magari un ricordo da
recuperare e su un'altra un minigioco come la discesa sullo scudo. In breve,
``riempirlo un po' di più''. Il docente conferma: gli spazi vuoti sono ``un
grosso problema degli ultimi giochi di Pokémon''. Daniele fa un confronto
interno: un'altra parte della mappa, il \textbf{deserto}, è invece costruita
\textbf{su più livelli}, ed è quindi uno spazio interessante; la montagna è
``un grande vuoto piano''.

\subsection{Francesca: i giochi psicologici e un labirinto di piani}

\textbf{Francesca} gioca ``un po' a tutto'', ma se deve scegliere preferisce
i giochi con un \textbf{valore psicologico} importante: molto narrativi, con
una psicologia di fondo da capire, non superficiale; horror indie, horror
psicologici. La sua forza è essere \textbf{estremamente inclusiva}: non giudica
gli altri e cerca di mettere tutti a proprio agio. La debolezza è
l'\textbf{introversione}, su cui sta lavorando. Come Daniele, non si sente
tagliata per fare la leader, pur sapendo comunicare e fare da mediatrice. Le
interessano sia la parte visiva e creativa sia quella tecnica e di
programmazione: per questo ha scelto il game design, ``perché più o meno
tocca tutti i punti''.

\paragraph{Il livello: Slime Rancher.} Il livello che non le è piaciuto per
niente è un'area del primo \textbf{\emph{Slime Rancher}}.\footnote{La
trascrizione riporta ``Dreamland e The Grey Lightning''; nel primo
\emph{Slime Rancher} esiste un'area chiamata \emph{The Grey Labyrinth}, che
corrisponde alla descrizione, ma il nome esatto pronunciato da Francesca non è
verificabile.} Non capisce il livello, e le ha fatto rigiocare la zona
``tantissime volte'': la mappa non è affatto chiara e, come nell'esempio del
docente, è tutto basato su \textbf{più piani intrecciati}, senza veri
collegamenti. Per passare da un punto all'altro o si trova l'unica strada
giusta ``tra i mille'' passaggi possibili, o ci si mette troppo tempo. Ed è
paradossale per un gioco di \emph{farming} ed esplorazione: va contro ciò che
il gioco stesso dovrebbe essere.

\subsection{Antonino: i roguelike e la tempesta di Frostpunk 2}

L'ultimo è \textbf{Antonino Etiopia} (``c'è sempre una persona che deve
essere ultima''). Gli andrebbe bene qualsiasi gioco, ma sceglierebbe i
\textbf{roguelike}, per le sinergie e le scelte; e, data la loro popolarità
nei giochi indie, avrebbe ``molte occasioni di provarci''. La sua forza è
lavorare in gruppo: rende meglio perché, se lavorasse male, la cosa ricadrebbe
sugli altri. ``Il senso di colpa è un eccellente motivatore.'' La debolezza,
che va un po' contro la forza, è l'introversione. Specializzazione: deve
ancora ambientarsi nel mondo del game dev.

\paragraph{Il livello: la città nella tempesta.} Il livello che gli ha fatto
abbandonare il gioco viene da \textbf{\emph{Frostpunk 2}}. Il docente ricorda
che \emph{Frostpunk} è una serie di \textbf{gestionali strategici} in cui si
governa una città in un'era glaciale, con risorse limitatissime e una
capacità di riscaldare gli ambienti molto ridotta: bisogna fare scelte
difficili su chi salvare, che cosa sacrificare, che cosa bruciare. Dopo che la
città sopravvive alla prima tempesta, bisogna andare a esplorare un'altra
città per recuperare risorse. Si hanno un certo numero di anni per
prepararsi, ma \textbf{non si ha idea di che cosa ci sia}, perché sopra la
città c'è una tempesta. Si arriva impreparati, e si perde tutto il lavoro.

\paragraph{L'analisi.} Che non si sappia che cosa si andrà ad affrontare è
``chiaramente intenzionale''. Ma qual è l'effetto? Ci si adatta, si resiste,
ci si prepara a perdere tutto e ad andare avanti comunque? O si ottiene
l'effetto ``perdo, ricomincio il gioco, e adesso che so cosa aspettarmi lo
faccio meglio''? Antonino aveva pensato la stessa cosa, ma il suo salvataggio
si era anche corrotto due volte, e ricominciare era diventato pesante (``e lì
c'è poco da fare: quando qualcosa si rompe, è difficile dire se è fatto
apposta'').

\paragraph{La proposta: un indizio.} Antonino darebbe, dopo la tempesta,
almeno \textbf{una dritta} su che cosa preparare. Il docente suggerisce di
inserirla nella \textbf{narrativa}: un sopravvissuto scappato da quella città
che racconta che cosa c'è (``ci sono dei dinosauri allergici allo zafferano:
manda la gente con i coltelli allo zafferano''; ``sparo cavolate, però è
possibile''). Qualcosa che permetta di \textbf{fare un piano} invece di
buttarsi nel vuoto totale. \emph{Frostpunk} è una serie hardcore, e mettere
bastoni fra le ruote è lo scopo del gioco; ma qui sembra ``un po' troppo
hardcore'', il tipo di cosa che ``devi saperlo prima: ci muori una volta e
riprovi''. Arrivare senza le strutture giuste significa sacrificare un sacco
di persone, e in un gioco così sbagliare una scelta tanto importante può
condannare l'intera partita: per quanto ricorda il docente, non ci sono
\textbf{meccaniche di recupero}. Se le cose vanno troppo male si entra in una
\textbf{spirale} in cui si continuano a perdere risorse senza che la partita
finisca.

% ---------------------------------------------------------------------
\section{Perché presentarsi: parlare in pubblico}
\label{gd8:sec:pubblico}

``Congratulazioni, avete passato la prima esercitazione del corso.'' Il giro
di presentazioni aveva due scopi. Il primo, dichiarato con ironia: il
docente fa fatica a ricordarsi i nomi e si dà una scusa per ripeterli (``Valerio'').
Il secondo è serio: \textbf{abituarsi a parlare in pubblico}.

Se si vuole fare qualunque cosa nei videogiochi, arriverà il momento in cui
bisognerà \textbf{presentare quello che si sta facendo}, e presentarlo a più
persone. Ci saranno molte occasioni per allenarsi, ma questa deve essere la
prima cosa che resta impressa: bisogna imparare a presentare sé stessi e il
proprio lavoro agli altri. Vale in quasi ogni settore lavorativo, ``ma
soprattutto nei settori creativi''.

\paragraph{La prima lezione seria.} ``Prima lezione seria di questo corso'':
\begin{quote}
\textbf{un designer mediocre, ma in grado di fare bei discorsi e di convincere
la gente, avrà sempre più successo nel proporre le proprie idee di un designer
geniale che non riesce a parlare in pubblico.}
\end{quote}
``Mi dispiace, la vita è così.'' Bisogna trovare un modo per superare le
proprie difficoltà nel parlare in pubblico: è ``assolutamente necessario''.
La buona notizia è che, a giudicare dal giro appena fatto, più o meno tutti
hanno una base di partenza; si tratta di migliorarla.

% ---------------------------------------------------------------------
\section{Che cos'è il level design}
\label{gd8:sec:cos-e}

\paragraph{La definizione.} ``Terra-terra'': il \textbf{level design} è la
disciplina che si occupa della \textbf{creazione di livelli, ambienti,
locali}, di qualunque tipo di spazio all'interno di un videogioco.

\paragraph{Il level editor.} Di solito lo si fa con un \textbf{level editor}.
``Grazie a Dio non sono più gli anni Ottanta'': oggi non si devono più fare i
livelli su carta a quadretti, prendendo le coordinate, convertendole e
inserendo tutto a mano. Per \emph{Super Mario Bros.}, per esempio, non c'era un
editor di livelli\footnote{La trascrizione riporta l'anno 1984; \emph{Super
Mario Bros.} è uscito nel 1985.}: c'era il codice. Ci si segnava le cose, si
scriveva il codice, e se un blocco finiva nel punto sbagliato bisognava
tornare nel codice e modificare il punto in cui il blocco era scritto.

\paragraph{Gli editor nei giochi.} Alcuni giochi includono un editor di
livelli che permette di crearne di nuovi all'interno del gioco stesso: sono
\textbf{ottimi per allenarsi} a fare level design con un set di meccaniche già
pronto.

\paragraph{Un lato artistico e un lato tecnico.} La cosa più importante da
ricordare è che il level design \textbf{ha un lato artistico e un lato
tecnico}. Lo si può trovare anche sotto altri nomi, come
\textbf{environment design} o \textbf{game mapping}: nel settore ``un sacco di
roba ha nomi diversi''.

\paragraph{Il blockout.} Gli screenshot della slide mostrano l'aspetto di un
videogioco \textbf{mentre viene sviluppato}: sono fasi di \textbf{blockout},
cioè fasi in cui non è ancora stata inserita l'arte finale e si usano forme
geometriche di base per cominciare a strutturare il livello. Sarà uno degli
argomenti del percorso.

\paragraph{Approfondimento: ``applied game design''.} Christopher Totten,
nell'introduzione di \emph{An Architectural Approach to Level Design} (il
libro che il docente consiglierà a fine lezione, sezione~\ref{gd8:sec:extra}),
raccoglie alcune definizioni classiche. Per Rudolf Kremers lo scopo del level
design è ``interpretare le regole del gioco e tradurle in un costrutto (un
livello) che faciliti al meglio il gioco'': il level design è \emph{game design
applicato}. Jay Wilbur e John Romero dicono che il level design è ``dove la
gomma tocca l'asfalto'', e che il level designer è ``in larga parte
responsabile dell'implementazione del gameplay''. Totten propone una
definizione più spaziale: il level design è ``l'esecuzione ragionata del
gameplay nello spazio di gioco, perché i giocatori lo abitino''. Ricorda anche
che molti designer hanno cominciato proprio come suggerisce il docente,
creando livelli personalizzati (\emph{mod}) con gli strumenti di giochi
esistenti.

% ---------------------------------------------------------------------
\section{Che cosa fa un level designer}
\label{gd8:sec:cosa-fa}

\subsection{La risposta breve}

``Un level designer si occupa di \textbf{creare spazi}.'' Spazi virtuali: in
2D, in 3D, ``in 1D, in 4D: esistono. Delirio totale, ma esistono''.

\subsection{La risposta lunga: un termine ombrello}

La risposta lunga è che ``level design'' è un \textbf{termine ombrello}: una
disciplina gigantesca che coinvolge diverse specializzazioni. A seconda del
contesto, chi lavora al level design può occuparsi di cose molto diverse fra
loro (figura~\ref{gd8:fig:ombrello}).

\begin{description}
  \item[Il blockout degli spazi.] ``Questo cubo è la casa, quest'altro cubo
        è il grattacielo, questo parallelepipedo più allungato è il ponte'':
        la prima struttura del livello, fatta di forme di base.
  \item[Le luci.] Illuminare non vuol dire ``prendo il sole, metto il sole,
        problema risolto'': le cose vanno illuminate \textbf{con cura}. Il
        level design si sovrappone molto alla \textbf{fotografia} e al
        \textbf{cinema}: nel cinema, le luci che illuminano una scena non
        sono mai quelle che si vedono. Ci sono sempre lampade fuori
        dall'inquadratura che illuminano con attenzione un personaggio o una
        zona, c'è la luce riflessa e soffusa che fa stagliare meglio i
        personaggi contro le ombre. Anche questo è un passaggio del level
        design: il \textbf{lighting pass}.
  \item[L'arredamento degli scenari.] Occuparsi di come
        \textbf{arredare} gli spazi, il cosiddetto \emph{set
        dressing}.\footnote{Il termine pronunciato dal docente è storpiato
        nella trascrizione (``la business plan''); dal contesto (``occuparsi di
        come arredare gli scenari'') si tratta con ogni probabilità del
        \emph{set dressing}, o di un \emph{dressing pass}.}
  \item[La documentazione.] Una cosa che fanno tutti i designer, e i level
        designer sono designer, è \textbf{scrivere documentazione}: da
        passare ad altri o da usare per tenere d'occhio il proprio lavoro.
        Layout, mappe, appunti, \textbf{diagrammi di flusso} che permettono
        di visualizzare la comunicazione tra le zone, la progressione nel
        livello, le azioni del giocatore. Un lavoro ``quasi burocratico''
        (si veda il capitolo~\ref{cap:brainstorming-flowchart} sui
        flowchart).
  \item[Il test.] Soprattutto, si \textbf{testa}. Le cose non si fanno ``su
        carta, poi in engine e via'': si testa, si fa \emph{playtest}, si
        prova un problema, si cambia, si riprova per vedere se un'altra
        opzione era migliore; si manda ai tester, e se c'è la possibilità si
        fa fare un giro al \textbf{QA}, persone esperte nel trovare
        \emph{bug}. È un processo \textbf{iterativo}: ``niente resta mai fisso
        nella prima forma che è stata creata''. Vale per le meccaniche come
        per i livelli.
  \item[La narrazione ambientale.] Vale anche per gli elementi più
        artistico-narrativi, come l'\textbf{environmental narrative}.
\end{description}

\paragraph{I criceti assassini.} L'esempio del docente sulla narrazione
ambientale merita di essere riportato. Nel livello c'è una registrazione:
``Oh no, i criceti sono scappati dalla gabbia e hanno sete di sangue. Hanno
già preso Frank, ma io non mi farò beccare. Se qualcuno sente questo
messaggio, il codice della porta che porta alle gabbie è 723. Li sento: stanno
graffiando le pareti. No\dots''. Chi scrive la registrazione non è
necessariamente il level designer; ma è il level designer che decide
\textbf{dove piazzarla}. Davanti alla porta 723? Oppure dall'altro lato della
mappa, così che il giocatore pensi: ``ah, 723 è il codice della porta, devo
tornare indietro, e attento che adesso ci sono i criceti assassini in giro''?
La stessa informazione, collocata in due punti diversi, produce due
esperienze diverse: una porta che si apre subito, oppure un ritorno carico di
tensione.

\begin{figure}[htbp]
\centering
\begin{tikzpicture}[
  font=\small,
  gd8core/.style={draw, very thick, ellipse, align=center, fill=blue!10,
                  inner sep=4pt},
  gd8leaf/.style={draw, thick, rounded corners=4pt, align=center,
                  text width=3.1cm, inner sep=4pt, fill=gray!6},
  gd8link/.style={thick}
]
  \node[gd8core] (ld) {\textbf{Level design}\\{\footnotesize creare spazi}};
  \node[gd8leaf, above left=1.0cm and 0.9cm of ld] (bo)
       {\textbf{Blockout}\\{\footnotesize forme di base: la casa,\\il grattacielo, il ponte}};
  \node[gd8leaf, above right=1.0cm and 0.9cm of ld] (lu)
       {\textbf{Luci}\\{\footnotesize il \emph{lighting pass},\\come nel cinema}};
  \node[gd8leaf, left=1.2cm of ld] (sd)
       {\textbf{Set dressing}\\{\footnotesize arredare gli scenari}};
  \node[gd8leaf, right=1.2cm of ld] (doc)
       {\textbf{Documentazione}\\{\footnotesize layout, mappe, appunti,\\diagrammi di flusso}};
  \node[gd8leaf, below left=1.0cm and 0.9cm of ld] (te)
       {\textbf{Test}\\{\footnotesize playtest, tester, QA:\\un processo iterativo}};
  \node[gd8leaf, below right=1.0cm and 0.9cm of ld] (en)
       {\textbf{Narrazione ambientale}\\{\footnotesize dove piazzare\\la registrazione}};
  \foreach \n in {bo,lu,sd,doc,te,en} \draw[gd8link] (ld) -- (\n);
\end{tikzpicture}
\caption{``Level design'' è un termine ombrello: a seconda del team e del
progetto, chi lo pratica può occuparsi di una o più di queste attività.}
\label{gd8:fig:ombrello}
\end{figure}

% ---------------------------------------------------------------------
\section{Lo strumento del level designer}
\label{gd8:sec:strumento}

\subsection{Il braccio armato dello sviluppo}

In un team di sviluppo, il level designer è quello che \textbf{mette di più
le mani nei giri}: per motivi pratici passa quasi tutto il tempo in engine, a
piazzare, sistemare, testare. ``Noi siamo il \textbf{braccio armato} dello
sviluppo, noi siamo quelli che zappano il terreno.'' La domanda, allora, è:
\textbf{qual è la nostra zappa?} Qual è lo strumento che un level designer
deve assolutamente conoscere per lavorare in un team?

\subsection{La parola magica: ``dipende''}

La prima risposta della classe è ovvia: l'\textbf{engine} in cui si crea il
gioco. Ma quale engine? \textbf{Dipende dall'azienda}, da quello che ha scelto
per creare il gioco. E qui il docente introduce ``la parola magica'' del corso:
nel corso dell'anno farà moltissime domande, e a moltissime di queste l'unica
vera risposta sarà \textbf{``dipende''}. ``Voglio che vi fissiate la parola
magica del mio corso in testa.''

L'engine potrebbe essere:
\begin{itemize}
  \item \textbf{Unreal};
  \item \textbf{Unity};
  \item \textbf{Godot}, ``che sta prendendo molto bene'';
  \item un \textbf{framework personalizzato}\footnote{La trascrizione riporta
        ``un custom framework fatto in 3GS''; potrebbe trattarsi di una
        libreria web come \emph{three.js}, ma non è verificabile.};
  \item \textbf{GameMaker}: \emph{Undertale} e \emph{Deltarune} sono fatti
        con GameMaker;
  \item un \textbf{engine proprietario}. Chi vuole lavorare in Capcom, per
        esempio, si troverà a lavorare con l'engine interno con cui sono
        stati fatti i \emph{Resident Evil} recenti e \emph{Monster
        Hunter}.\footnote{La trascrizione dice ``è probabile che si trovino a
        lavorare con l'Unreal Engine'', ma subito dopo parla dell'engine
        ``con cui è stato fatto [Resident Evil] e con cui è stato fatto Monster
        Hunter''. L'engine interno di Capcom è il RE Engine; è probabile che
        la prima frase sia un errore di trascrizione per ``RE Engine''.}
\end{itemize}
Lo strumento indispensabile, quindi, \textbf{non può essere un engine
preciso}, perché non si sa mai quale si andrà a usare. Nel corso si userà
\textbf{Unreal}, perché è considerato lo standard dell'industria ``per
adesso''; ma questo non significa che sia indispensabile conoscerlo. Ciò che
si studia su Unreal ha applicabilità anche negli altri engine, magari con nomi
diversi e procedimenti diversi.

\subsection{La risposta vera: il level designer}

``L'unico strumento di cui un level designer ha bisogno per fare il suo lavoro
è \textbf{il level designer}.''

Ecco lo scopo delle domande di inizio lezione: non servivano solo a conoscere
la classe o a capire se fossero ``tutti fan di rhythm game e di simulatori di
gioco d'azzardo''. Servivano a far capire che \textbf{lo strumento principale}
che si userà in tutta la produzione di videogiochi \textbf{siete voi}. Conoscere
le proprie forze e le proprie debolezze serve ad applicarle a contesti spesso
molto diversi tra loro, perché \textbf{sono il contesto e le necessità del
progetto a decidere quali strumenti si usano}.

\paragraph{Imparare a imparare.} ``La vostra capacità di imparare è il vostro
asset migliore.'' Il docente può insegnare per tre anni come si usa Unreal;
poi si va a lavorare in Capcom, o in FromSoftware, che usa un suo engine
sviluppato dagli anni Duemila, totalmente diverso da Unreal. La capacità da
sviluppare in questi tre anni non è saper usare Unreal: è \textbf{sapere come
si impara}. La caratteristica principale dello sviluppatore di videogiochi è
essere \textbf{uno studente continuo}: la flessibilità mentale di imparare in
fretta strumenti diversi, saper fare ricerca, saper fare pratica, tenersi
allenati, tenere d'occhio i cambiamenti.

\paragraph{Due esempi.} Epic ha annunciato che molto cambierà con la
\textbf{versione 6} di Unreal, ``quando o se esce'': quando succederà, se si
deciderà di passarci, bisognerà attivarsi. E un gruppo di studenti del terzo
anno ha chiesto di fare il proprio progetto in \textbf{Unity}: metà del gruppo
ci si trovava meglio, l'altra metà non l'aveva mai toccato. Quest'altra metà
``si è rimboccata le maniche'' ed è andata a impararlo.

``Imparate a imparare: questa è la cosa più importante che dovete saper
fare.'' (Seguito da un'autoironia: ``vedo che alcuni di voi hanno la
lacrimuccia\dots{} no, siamo tutti cinici, col cuore raggrinzito come una
prugna.'')

% ---------------------------------------------------------------------
\section{Come si impara il level design}
\label{gd8:sec:imparare}

A livello pratico, che cosa bisogna fare per imparare il level design? Il
docente indica alcune aree su cui concentrarsi.

\subsection{La teoria}

Il vantaggio della teoria è che \textbf{si applica a qualunque cosa}. I
principi di level design valgono che si lavori in Unity, in Unreal o in Godot,
che si stia facendo un FPS o un gioco di corse.

\paragraph{Il landmark.} L'esempio è il \textbf{landmark}: ``un landmark è un
landmark è un landmark'', ha sempre la stessa funzione in qualunque tipo di
gioco. Un landmark è \textbf{un elemento dello scenario distinguibile e
visibile dalla distanza, che aiuta il giocatore a orientarsi}. La
\textbf{Madonnina} è il landmark di Milano. Lo si vedrà con attenzione nelle
lezioni future; ma si noti che è esattamente ciò che mancava nella Cattedrale
del Sangue Sacro.

\paragraph{Rifiutare una regola, sapendolo.} Anche quando si lavora a un gioco
che \textbf{rifiuta intenzionalmente} queste regole, bisogna sapere che le si
sta rifiutando. Un horror alla \emph{Silent Hill}, con la nebbia che impedisce
di vedere a distanza, lo fa apposta: non è che i suoi autori non sappiano che
cosa sia un landmark. C'è una conoscenza di questi elementi e un
\textbf{rifiuto di usarli}, che fa parte della loro intenzione.

\paragraph{Approfondimento: Lynch e l'immagine della città.} Totten dedica
al tema una parte del suo libro, ispirata a Kevin Lynch e al suo studio
sull'\emph{immagine della città}: gli abitanti si orientano grazie a pochi
elementi riconoscibili, tra cui i \emph{percorsi} e i \emph{landmark}, punti di
riferimento visibili da lontano. Nel capitolo sul \emph{pathfinding} con le
``architectural weenies'' (l'espressione con cui Walt Disney indicava gli
elementi visibili da lontano che attirano il visitatore, come il castello al
centro del parco), il landmark non serve solo a orientarsi ma anche a
\textbf{attirare} il giocatore verso una meta.

\subsection{La pratica}

``La teoria non vale niente senza la pratica.'' Bisogna mettere le mani in
engine. C'è chi è più interessato alla teoria; ma finché si seguono le lezioni
di questo corso si fa level design, almeno parzialmente, almeno per imparare.
``Più pratica fate, più imparate a fare quello che dovete fare, più diventate
veloci, più cose potete fare.'' E nei videogiochi ``la velocità è tutto''.

\subsection{L'expertise}

L'\textbf{expertise} non è solo esperienza: è \textbf{l'insieme di tutto
quello che si sa fare}. Chi sa fare gli \emph{shader} può farsi un sacco di
cose da solo, in fretta, senza aspettare altri. Chi sa fare i suoni non deve
passare ore a cercare materiale, o rassegnarsi a risultati molto peggiori di
quelli che otterrebbe un esperto.

\subsection{L'esperienza sul campo}

``Sembra stupido'', ma vagare, esplorare, fare \emph{trekking}, girare per le
città, osservare l'ambiente rende un level designer migliore. \emph{The Legend
of Zelda} è nato perché \textbf{Shigeru Miyamoto} voleva ricreare
l'esperienza che aveva avuto da bambino, esplorando la campagna giapponese
intorno a casa sua.\footnote{La trascrizione riporta ``Minamoto'' e
``Zelda di terraferma''; l'aneddoto è quello, noto, di Miyamoto e delle sue
esplorazioni d'infanzia nei dintorni di Kyoto.}

\subsection{Il gruppo}

Tutto questo (esperienza, ricerca, studio) non serve a niente se non lo si
mette in pratica; e la pratica, in questo corso, avverrà spesso \textbf{in
gruppo}. C'è chi vuole lavorare da solo, ma finché si segue il corso si lavora
anche con gli altri, ed è fondamentale saperlo fare. ``Vi assicuro che
finirete questi tre anni che qualcuno di voi odierà qualcun altro.
Statisticamente è impossibile che non succeda.'' Ma anche lavorare con gli
altri richiede pratica: è la pratica che rende veloci, che permette di
ottenere risultati nei tempi previsti e di conoscere gli strumenti.

\subsection{La parola che si ripete}

``Siete persone intelligenti: starete notando che sto ripetendo molto spesso
una certa parola.'' La parola è \textbf{pratica}. L'ultima cosa, ma non ultima
per importanza: la cosa più importante quando si fa level design è
\textbf{fare pratica}. Più si mettono le mani in pasta, meglio si lavora.
``Tutto il resto dipende.''

% ---------------------------------------------------------------------
\section{Due curve da conoscere in anticipo}
\label{gd8:sec:curve}

Il docente mostra due grafici tratti dalla ricerca sociologica e
psicologica, che descrivono ciò che succederà a quasi tutti nel corso di
studi. Li presenta adesso perché ``è inutile affrontarlo quando succede'':
prima ci si prepara, meglio è.

\subsection{Il picco della fiducia}

Si parte tutti più o meno da zero, o comunque da un livello molto basso. Nel
corso del primo anno si impara una serie di cose: come mettere le mani in
pasta, come analizzare i giochi. Più si impara, più cresce la
\textbf{fiducia} di sapere che cosa si sta facendo, più ci si sente esperti:
``più vi sentirete fighi, più andrete dai vostri amici su Discord a dire: eh,
quel livello non è fatto troppo bene, io avrei messo un altro falò''.

Poi, di solito tra la fine del primo anno e l'inizio del secondo, si arriva al
\textbf{picco}, e ci si scontra con la realtà: non si è ``così fighi'' come si
credeva. Ci si rende conto di quanto manca ancora da imparare, di quanto la
propria conoscenza sia superficiale, di quanto i risultati siano scarsi
rispetto alle speranze. ``Alessandro lo sa.'' \textbf{Più si sa, più ci si
rende conto di quanto poco si sa.} È \textbf{fisiologico}: il docente si
aspetta che ognuno abbia ``una botta di energia, figaggine, quanto sono bravo,
quanto sono bello'', e poi, tra primo e secondo anno, ``la depressione''.
``Non c'è niente di male, succede.'' Se si riesce a superare quel picco, si
possono fare grandi cose.\footnote{Il grafico della slide non è disponibile e
la trascrizione non ne riporta il titolo esatto (``il picco della realtà'');
la descrizione è ricostruita dalle parole del docente.}

\subsection{L'effetto Dunning-Kruger}

Il secondo grafico è l'\textbf{effetto Dunning-Kruger}, anch'esso un effetto
osservato: la \textbf{fiducia} di una persona in quanto sa di un argomento
segue una curva per cui, \textbf{meno si sa, più ci si sente sicuri} nel
parlarne (figura~\ref{gd8:fig:dk}).

\paragraph{L'allenatore di calcio.} L'esempio del docente: parenti o amici che
guardano il calcio e parlano con enorme sicurezza di quanto sia scemo
l'allenatore, di quale giocatore avrebbe dovuto sostituire, ``e non capisco
perché''. ``Quello è l'effetto Dunning-Kruger in azione.'' Ne sanno di più di
una persona che fa calcio da trent'anni? Probabilmente no.

\paragraph{Il giocatore che non ha mai fatto un gioco.} Appena si impara poco
di un argomento, la fiducia (``io so quasi tutto di questa roba, sono un
esperto'') schizza in alto. La gente che gioca ai videogiochi ma non li ha mai
fatti si trova più o meno lì: ``fiduciosissima'' nel dire ``questo gioco fa
schifo, avrebbero dovuto cambiare questa cosa, avrebbero dovuto portare i
giocatori di una mappa da 12 a 128''. È chiaro che non sanno di che cosa
parlano. Poi, più si impara, più ci si rende conto di quanto sia limitata la
propria conoscenza: il punto più basso della curva è quando si sono imparate
molte cose e si capisce quanto ne manchino ancora.

\paragraph{L'obiettivo.} ``Il vostro obiettivo all'interno di questo corso di
studi è \textbf{sopravvivere ai due punti centrali di questo grafico, senza
abbandonare la nave}.'' E, aggiunge, ``non rendetemi la vita impossibile
finché siete ancora nella prima metà del grafico'': durante l'anno ci si
sentirà ``fighissimi'', e lui sarà lì a dire ``bravi, bravi'', aspettando che
la prima fase passi. ``L'annuncio di servizio ve l'ho fatto; il resto sta a
voi.''

\begin{figure}[htbp]
\centering
\begin{tikzpicture}[
  font=\small,
  gd8ax/.style={-{Stealth[length=2.5mm]}, thick},
  gd8lab/.style={font=\footnotesize, align=center},
  gd8dot/.style={circle, fill=black, inner sep=1.4pt}
]
  \draw[gd8ax] (0,0) -- (11,0) node[below left, gd8lab] {competenza, esperienza};
  \draw[gd8ax] (0,0) -- (0,5.3) node[left, gd8lab] {fiducia};
  \draw[very thick, blue!60!black]
    (0.1,0.3) .. controls (0.9,5.6) and (1.9,5.4) .. (2.8,2.6)
              .. controls (3.4,0.8) and (4.4,0.6) .. (5.4,1.2)
              .. controls (7.0,2.2) and (8.8,3.3) .. (10.6,3.8);
  \node[gd8dot] at (1.45,4.25) {};
  \node[gd8lab, anchor=west] at (2.05,4.75)
       {``so quasi tutto'': il tifoso, il giocatore\\che non ha mai fatto un gioco};
  \node[gd8dot] at (4.5,0.85) {};
  \node[gd8lab, anchor=north] at (4.5,-0.15)
       {``quanto poco so'':\\il punto più basso};
  \node[gd8lab, anchor=south] at (8.6,3.45)
       {la risalita: la fiducia\\torna, fondata sulla pratica};
  \node[gd8lab, text=red!70!black, anchor=west] (gd8fy) at (4.0,3.3)
    {fine primo anno /\\inizio secondo anno};
  \draw[thick, red!70!black, -{Stealth[length=2mm]}] (gd8fy.west) -- (2.75,2.85);
\end{tikzpicture}
\caption{Schema dell'effetto Dunning-Kruger come lo descrive il docente: con
poche conoscenze la fiducia schizza in alto; imparando si precipita nel punto
in cui ci si rende conto di quanto poco si sa; solo dopo risale. L'obiettivo
del corso è superare la parte centrale della curva ``senza abbandonare la
nave''. Le proporzioni sono indicative.}
\label{gd8:fig:dk}
\end{figure}

\paragraph{Approfondimento: che cosa dice davvero lo studio.} L'effetto prende
il nome da David Dunning e Justin Kruger, che nel 1999 pubblicarono uno studio
in cui le persone meno competenti in un compito tendevano a
\textbf{sovrastimare} la propria prestazione, e le più competenti a
sottostimarla leggermente. La curva ``a montagna e valle'' che circola sotto
questo nome (e che si vede anche nella slide) è una divulgazione molto
diffusa, più che il grafico originale dello studio; ma rende bene l'idea che
il docente vuole trasmettere: la sicurezza non è una misura della competenza.

% ---------------------------------------------------------------------
\section{Tre corollari}
\label{gd8:sec:corollari}

\subsection{Primo: non esistono regole fisse}

Legato alla parola magica: \textbf{nel design non esistono regole fisse}, che
si applichino invariabilmente a ogni cosa. Lo scopo non è mai trovare la
risposta perfetta da applicare a ogni situazione: è \textbf{isolare la domanda
e scoprire la ragione} per cui si sta facendo una scelta. Quella ragione la si
può poi declinare a seconda della situazione. Progetti diversi, dimensioni
diverse, disponibilità economiche o di risorse diverse, persone diverse nel
team portano a risposte diverse.

\paragraph{L'artista del fantasy.} Se l'artista che crea i modelli 3D non ha
mai fatto fantascienza, fare un gioco di fantascienza ``magari non è il
massimo''. Conviene considerare un'ambientazione fantasy, ``perché è quello che
il nostro disegnatore disegna da quando aveva cinque anni''.

\subsection{Secondo: il design è una disciplina olistica}

``Chi conosce questa parola?'' Silenzio. L'\textbf{olismo} è il considerare
qualcosa non come la somma di più parti separate, ma come \textbf{un unico
insieme}, diverso dalla somma delle parti. Un designer deve capire l'insieme,
non solo le parti.

L'esempio pratico, dal level design: bisogna sapere
\begin{itemize}
  \item creare una \textbf{conformazione del terreno} interessante;
  \item costruire una \textbf{casa} che abbia, almeno visivamente, un senso
        architettonico, con gli strumenti a disposizione;
  \item posizionare gli \textbf{elementi sonori}: se c'è un fiume, il suono
        deve accompagnare il giocatore mentre si sposta lungo il fiume;
  \item posizionare l'\textbf{illuminazione}: la luce di una lampada non sta
        ``dentro la lampada'', va regolata con attenzione, e certi oggetti
        vanno esclusi dal proiettare ombre, ``altrimenti diventa proprio
        brutto vedere tutte le ombre sparate sui muri'';
  \item costruire \textbf{puzzle}, posizionare \textbf{nemici}.
\end{itemize}
Tutte queste cose formano un insieme che è più della somma delle parti. Un
level designer deve avere almeno una \textbf{conoscenza di base di tutti gli
strumenti e di tutte le discipline} con cui lavora; e più sa delle altre
discipline, meglio riesce a \textbf{comunicare con i colleghi} che se ne
occupano.

\subsection{Terzo: il design non è arte, è artigianato}

Qual è la differenza? L'\textbf{arte} è qualcosa che non può essere giudicato
pienamente secondo canoni puramente tecnici. Non si può dire che i
\emph{Sex Pistols} facessero schifo perché non sapevano fare
\emph{shredding} velocissimo come i \emph{Metallica}: sono due cose diverse,
c'è un intento artistico.

L'\textbf{artigianato} è diverso. Non si può dire ``questa sedia non è fatta
per sedercisi, perché se ti siedi si spacca, però guarda com'è bella e
artistica''. No: è una sedia, il suo scopo è che ci si possa sedere sopra. Se
non ci si può sedere, non si sta facendo una sedia: si sta facendo un'opera
d'arte, e allora non la si mette sotto il tavolo, la si mette su un
piedistallo. Nel design \textbf{si fa artigianato}: ciò che si costruisce ha
uno \textbf{scopo}, e richiede capacità ed esperienza artigianali.

\subsection{Sporcarsi le mani}

Da qui i consigli pratici:
\begin{itemize}
  \item \textbf{fare giochi}: più giochi si fanno, meglio si diventa come
        designer. Le \emph{game jam}, se si vuole, ma anche progetti un po'
        più lunghi;
  \item \textbf{ricreare} ciò che si vede, con gli strumenti che si
        studieranno nel corso;
  \item \textbf{trovare altre persone}, ``creare una scena'': andare alle
        serate di \emph{Milano Indies}, nei caffè in cui la gente si trova a
        programmare, costruirsi attorno un gruppo di persone che fanno la
        stessa cosa, ``perché si migliora sempre insieme''. I punk non
        stavano da soli in cameretta: si trovavano nei bar, andavano ai
        concerti insieme, si incontravano al parco, e hanno creato una delle
        scene musicali più importanti della storia;
  \item \textbf{farsi ispirare}: variare i giochi a cui si gioca, i film
        che si guardano, la musica che si ascolta.
\end{itemize}

% ---------------------------------------------------------------------
\section{La struttura del corso}
\label{gd8:sec:struttura}

Prima della parte pratica, un riassunto di come funzionerà il corso.
\begin{description}
  \item[Pratica in engine.] Si lavorerà molto in \textbf{Unreal Engine 5}.
        Ci saranno momenti di \textbf{pratica assistita}, in cui il docente
        mostra e guida, e momenti di \textbf{pratica a casa}.
  \item[I task.] Il docente assegna dei compiti, da consegnare entro una o
        due settimane a seconda del task, e poi li si rivede insieme in
        classe per correggere eventuali problemi.
  \item[La teoria.] Ci saranno momenti di studio teorico: \textbf{teoria} e
        \textbf{metodologia} del level design. Anche la teoria avrà momenti
        di pratica assistita, esercizi in classe ed esercizi a casa.
  \item[Le letture.] Qualche libro per chi vuole approfondire, non
        necessariamente obbligatorio: il docente tende a lasciare
        suggerimenti più precisi a chi decide di specializzarsi nel level
        design, e lo si deciderà durante l'anno in base alle preferenze.
\end{description}

\paragraph{La pausa.} Si è a metà lezione, ed è ora della pausa. Il docente ha
una ``malattia un po' pesante'' e mal di gola, quindi la fa un po' più lunga,
una ventina di minuti, ``sennò arrivo stasera che sputo sangue''. Al rientro
si passa a Unreal: bisogna averlo installato, in una versione recente (``dopo
la 5.5, 5.6 dovrebbe essere più o meno la stessa cosa''; al massimo cambia
leggermente l'interfaccia).\footnote{La registrazione della pausa contiene
quasi solo frammenti incomprensibili e conversazioni private degli studenti,
che non sono riportate.}

% ---------------------------------------------------------------------
\section{Il progetto in Unreal Engine}
\label{gd8:sec:progetto}

\subsection{Lanciare l'engine in anticipo}

Prima ancora di riprendere, il docente chiede di \textbf{lanciare Unreal}
negli ultimi minuti della pausa: ``ha la tendenza a metterci un pochino ad
aprirsi''. Va bene qualsiasi versione dalla 5.6 in su; sul computer dell'aula
non è installata la 5.8, quindi il docente usa la 5.6. (``Non dirmi che si sta
aggiornando. Ti prego.'')

\subsection{Creare il progetto}

Qualunque progetto va bene per ciò che si farà oggi. I passaggi:
\begin{enumerate}
  \item dal launcher si avvia Unreal Engine (5.6 o successivo): si apre il
        \textbf{Project Browser};
  \item si sceglie la categoria \textbf{Games} e il template \textbf{Third
        Person};
  \item si dà un nome al progetto (il docente usa ``primo anno'');
  \item si fa \textbf{Create} e si aspetta che il progetto venga creato.
\end{enumerate}

\paragraph{Lo starter content.} La slide mostra una versione più vecchia del
launcher, in cui lo \textbf{starter content} si includeva con una spunta al
momento della creazione. Nelle versioni nuove non c'è più quell'opzione, ma lo
si può aggiungere dopo (lo si vedrà tra poco, sezione~\ref{gd8:sec:content}).
Per oggi non serve.

\paragraph{Il computer che crolla.} Come nella lezione di programmazione
(capitolo~\ref{gd7:cap:blueprint}), il computer dell'aula non regge: l'editor
si blocca, il disco ``sta facendo un casino terrificante'', e il docente deve
riaprire tutto più volte, fino ad aprire un progetto qualsiasi già presente
(``incrociamo le dita''). Nel frattempo fa l'appello.

\paragraph{Le slide.} Le slide della lezione sono nel \textbf{drive
condiviso} degli studenti: cartella di Milano, primo anno, \emph{Lezioni},
\emph{Level Design}. Il docente le scrive in modo molto esteso proprio per
usarle anche come appunti, così che si possa rivedere tutto; ma consiglia
comunque di \textbf{prendere appunti}, ``perché è il modo migliore di fissarsi
le cose in mente''.

% ---------------------------------------------------------------------
\section{L'editor: muoversi e piazzare oggetti}
\label{gd8:sec:editor}

L'editor è diviso in \textbf{finestre}. Si possono riarrangiare e spostare
dove si preferisce, ma alcune vanno conosciute subito
(si vedano la figura~\ref{ue1:fig:editor} e la
sezione~\ref{ue1:sec:interfaccia}).

\subsection{La viewport}

La finestra più importante è quella \textbf{centrale}, che mostra l'aspetto
del livello: la \textbf{viewport}, ``il cuore del nostro editor''. È
probabilmente l'unica finestra che resterà sempre aperta, comunque si decida di
organizzare l'editor. Cliccando sugli oggetti li si \textbf{seleziona}.

\paragraph{Che cos'è un actor.} Un \textbf{actor} è \textbf{un oggetto
all'interno del gioco}, nel livello. La definizione verrà precisata tra poco,
distinguendolo dall'\emph{asset} (sezione~\ref{gd8:sec:content}).

\paragraph{Muoversi.} Per spostarsi nella viewport:
\begin{itemize}
  \item si \textbf{tiene premuto il tasto destro} del mouse e si muove il
        mouse per \textbf{ruotare la telecamera};
  \item tenendo premuto il tasto destro, con \textbf{W, A, S, D} ci si sposta
        ``come in un first-person shooter'';
  \item sempre con il tasto destro premuto, \textbf{Q} ed \textbf{E} spostano
        la telecamera in \textbf{verticale}.
\end{itemize}
``Prendeteci la mano'': sono i controlli di base per esplorare la scena che
Unreal crea automaticamente per ogni nuovo progetto.

\subsection{Place Actors e le forme di base}

A sinistra c'è la finestra \textbf{Place Actors}, che permette di selezionare
degli actor da inserire nel livello. ``Oggetti'' non vuol dire necessariamente
oggetti fisici: possono essere \textbf{luci}, \textbf{suoni},
\textbf{volumi} (che si vedranno più avanti). Si afferrano, si trascinano e si
piazzano nel mondo di gioco.

Gli actor sono divisi in \textbf{categorie} (nelle versioni recenti l'elenco
delle categorie è verticale; gli screenshot della slide vengono da una
versione un po' precedente, ma le funzionalità sono le stesse). Se lo schermo è
piccolo, \textbf{tre puntini} indicano che ci sono altre categorie nascoste.
Quella che interessa oggi è \textbf{Shapes} (sotto \emph{More}, nelle versioni
recenti), che contiene le \textbf{forme di base}.

Si prende il \textbf{cubo}, lo si trascina in uno spazio libero della scena:
``congratulazioni, avete fatto lo stesso passo che fa un bambino di due anni
quando gioca con i cubi''. Un passo alla volta.

\subsection{Il gizmo}

Quando si trascina un oggetto nella scena compaiono tre frecce: un piccolo
\textbf{gizmo} che permette di spostarlo (figura~\ref{gd8:fig:gizmo}).
\begin{itemize}
  \item Afferrando una \textbf{freccia} si sposta l'oggetto lungo un solo
        asse: \textbf{X} (rossa), \textbf{Y} (verde) o \textbf{Z} (blu, la
        verticale).
  \item Afferrando il \textbf{centro}, lo si sposta in relazione alla
        telecamera, sul piano che la telecamera sta guardando.
  \item Afferrando uno dei \textbf{quadratini} tra due frecce, lo si sposta
        su \textbf{due assi contemporaneamente}: si evidenziano entrambe le
        frecce. Con il quadratino tra X e Y, per esempio, lo si sposta solo
        sul piano orizzontale, senza preoccuparsi di alzarlo o abbassarlo:
        se lo si era appoggiato a terra, resterà a terra.
\end{itemize}
``Prendeteci la mano, perché sarà la cosa che farete di più in questo
engine.'' Ogni volta che si seleziona qualcosa, anche un elemento dello
scenario, compare un gizmo, e lo si può spostare a piacere.

\begin{figure}[htbp]
\centering
\begin{tikzpicture}[
  font=\small,
  gd8arr/.style={-{Stealth[length=3mm]}, line width=1.6pt},
  gd8lab/.style={font=\footnotesize, align=left}
]
  \coordinate (o) at (0,0);
  \draw[gd8arr, red!80!black] (o) -- (2.6,-0.9) node[right] {\textbf{X}};
  \draw[gd8arr, green!50!black] (o) -- (2.4,0.9) node[right] {\textbf{Y}};
  \draw[gd8arr, blue!80!black] (o) -- (0,2.6) node[above] {\textbf{Z}};
  % quadratini dei piani
  \fill[yellow!70!orange, opacity=0.8] (0.75,-0.26) -- (1.45,0) -- (0.7,0.26) -- (0,0) -- cycle;
  \fill[cyan!40, opacity=0.8] (0,0.7) -- (0.7,0.96) -- (0.7,0.26) -- (0,0) -- cycle;
  \fill[magenta!35, opacity=0.8] (0,0.7) -- (0.75,0.44) -- (0.75,-0.26) -- (0,0) -- cycle;
  \fill[white, draw=black, thick] (o) circle (0.13);
  \node[gd8lab, anchor=west] at (3.6,2.1)
    {\textbf{frecce}: un solo asse\\(X rosso, Y verde, Z blu)};
  \node[gd8lab, anchor=west] at (3.6,0.9)
    {\textbf{quadratini}: due assi insieme\\(es. X--Y: resta sul pavimento)};
  \node[gd8lab, anchor=west] at (3.6,-0.3)
    {\textbf{centro}: sul piano\\perpendicolare alla telecamera};
  \node[gd8lab, anchor=west] at (3.6,-1.5)
    {\textbf{W} sposta \quad \textbf{E} ruota \quad \textbf{R} scala};
\end{tikzpicture}
\caption{Il gizmo di traslazione, schematizzato. Il punto in cui compare (il
\emph{pivot}) non è necessariamente il centro geometrico dell'oggetto: lo
decide chi ha creato la mesh.}
\label{gd8:fig:gizmo}
\end{figure}

\paragraph{Il pivot non è sempre al centro.} Cliccando in giro si nota subito
che non tutti i gizmo stanno al centro dell'oggetto: certe pedane della scena,
per esempio, lo hanno sulla parete frontale. Il punto in cui viene piazzato il
gizmo è \textbf{arbitrario}: lo decide chi fa la \emph{mesh}, e può essere in
qualunque punto. Lo si può anche spostare, se serve cambiare l'origine di una
mesh, ma se ne parlerà più avanti.

% ---------------------------------------------------------------------
\section{L'Outliner: tenere in ordine il livello}
\label{gd8:sec:outliner}

In alto a destra c'è l'\textbf{Outliner}: \textbf{la lista di tutti gli actor
presenti nel livello}. (Il docente usa la terminologia inglese, ``actor''
piuttosto che ``attori'': ``è più facile trovare risorse e abituarsi''.)

\paragraph{Ogni oggetto ha un nome.} Il cubo appena trascinato, nel caso del
docente, si chiama \emph{Cube2}, perché c'era già un altro cubo. Cliccandoci
sopra nell'Outliner lo si seleziona, e lo si ritrova nella viewport.

\paragraph{Le cartelle.} Tutto ciò che si trascina nel livello entra
nell'Outliner secondo una \textbf{gerarchia}: gli actor si possono spostare,
ordinare, raccogliere in \textbf{cartelle} (nella scena di esempio ce n'è una
chiamata \emph{Playground}). Una nuova cartella si crea con tasto destro,
\textbf{Create Folder} (il docente la chiama ``Bubba'', ``non so come
chiamarla''). Organizzare le cose come file in cartelle, come sul proprio
computer, aiuta a \textbf{tenere organizzato il livello}: ogni volta che si
mette qualcosa nel gioco, fare attenzione al \textbf{nome} e alla
\textbf{cartella}.

\paragraph{L'ordine alfabetico.} Nota importante: l'Outliner è
\textbf{ordinato alfabeticamente}. Non gli interessa che cosa è stato piazzato
prima o dopo, né dove: l'ordine è sempre alfabetico. Le cartelle servono a
raggruppare, e l'ordine alfabetico è quello in cui le cose verranno mostrate.

\paragraph{Selezionare più oggetti.} L'Outliner rispetta le regole di
Windows:
\begin{itemize}
  \item cliccando sul primo elemento e poi, con \textbf{Shift}, sull'ultimo,
        si seleziona tutto ciò che sta in mezzo (il docente seleziona così
        tutti gli oggetti della cartella \emph{Playground});
  \item con tasto destro su una cartella, \textbf{Select} $\rightarrow$
        \textbf{Immediate Children}, si selezionano rapidamente tutti gli
        elementi contenuti.
\end{itemize}

\paragraph{Rinominare.} Come in Windows, si seleziona l'oggetto e si preme
\textbf{F2} (il docente lo chiama ``Cubo Super Ultra''). In alternativa, un
secondo clic lento sul nome, oppure tasto destro, \emph{Rename}. Ma il
consiglio ``spassionato'' è imparare le \textbf{scorciatoie}: ``vi velocizzano
tanto il lavoro''. I tasti \emph{funzione}, quelli più in alto sulla
tastiera, vale la pena impararli.

\paragraph{Ctrl+Z.} Come in Windows, \textbf{Ctrl+Z} annulla l'ultima azione:
``abituatevi a usarlo, è la scorciatoia più utile''.

% ---------------------------------------------------------------------
\section{Il Content Browser: asset e actor}
\label{gd8:sec:content}

In basso c'è il \textbf{Content Browser}: \textbf{tutti gli asset che fanno
parte del gioco}. Qui il docente introduce una distinzione importante:
\begin{description}
  \item[Asset] un oggetto \textbf{salvato tra i file del progetto}, su disco;
  \item[Actor] un oggetto \textbf{nell'editor}, all'interno di un livello.
\end{description}
Gli asset stanno nel progetto e possono essere trascinati dal Content Browser
dentro un livello; lì diventano actor.

\paragraph{Aggiungere lo starter content.} Nelle versioni recenti lo starter
content non viene inserito automaticamente. Per aggiungerlo (passaggio che non
è nelle slide):
\begin{enumerate}
  \item tasto destro nel Content Browser;
  \item \textbf{Add Feature or Content Pack};
  \item nella sezione \textbf{Content}, \textbf{Starter Content}.
\end{enumerate}
È un insieme di texture, oggetti e altro materiale, fornito gratuitamente da
Epic con Unreal Engine, per cominciare a fare i propri progetti. Oggi non
serve, perché si lavora solo con le forme di base; ce ne si preoccuperà dalla
prossima lezione.\footnote{Il docente aggiunge un commento sulle scelte
recenti di Epic (``credo che tutti quelli che lavorano nel mondo dei
videogiochi si stiano lamentando delle scelte che sta facendo Epic in questi
ultimi mesi''), senza specificare a quali si riferisca.}

\paragraph{Le finestre si ridimensionano.} Come tutte le finestre, anche il
Content Browser si può allargare o stringere a piacere.

\subsection{Istanze e stampi}

Il docente ci torna sopra, perché è una distinzione da avere chiara. Un
\textbf{asset} è la versione salvata su disco di un oggetto del progetto. Un
\textbf{actor}, in termini tecnici, è una \textbf{instance} (istanza): una
\textbf{copia} di quella cosa salvata su disco, inserita nel livello. Il cubo
``Cubo Super Ultra'' è un'istanza del cubo; l'altro cubo nella scena è
un'altra istanza dello stesso cubo.

\paragraph{Le modifiche restano sull'istanza.} Quando si modifica un oggetto
nei \emph{Details} o con il gizmo, \textbf{si modifica l'istanza nel livello,
non l'asset nel progetto}. Se si allarga una sedia, le si cambia il materiale
e poi si trascina di nuovo la stessa sedia dal Content Browser, la nuova sarà
identica alla sedia di partenza. Le modifiche fatte nell'editor restano solo
su quell'istanza, e non vengono sovrascritte sull'originale.

L'immagine suggerita: pensare agli asset come a uno \textbf{stampo}, e alle
istanze nel livello come alle \textbf{copie} create da quello stampo
(figura~\ref{gd8:fig:stampo}).

\begin{figure}[htbp]
\centering
\begin{tikzpicture}[
  font=\small,
  gd8asset/.style={draw, very thick, rounded corners=3pt, align=center,
                   fill=orange!15, inner sep=5pt},
  gd8inst/.style={draw, thick, align=center, fill=gray!8, inner sep=4pt},
  gd8arr/.style={-{Stealth[length=2.5mm]}, thick},
  gd8lab/.style={font=\footnotesize\itshape, align=center}
]
  \node[gd8asset] (a) {\textbf{Asset}: \emph{Chair}\\{\footnotesize nel Content Browser}\\{\footnotesize (su disco, nel progetto)}};
  \node[gd8inst, right=2.4cm of a, yshift=1.3cm] (i1) {\textbf{Istanza 1}\\{\footnotesize scala 1, materiale originale}};
  \node[gd8inst, right=2.4cm of a] (i2) {\textbf{Istanza 2}\\{\footnotesize allargata, materiale cambiato}};
  \node[gd8inst, right=2.4cm of a, yshift=-1.3cm] (i3) {\textbf{Istanza 3}\\{\footnotesize trascinata dopo: identica all'asset}};
  \draw[gd8arr] (a.east) -- (i1.west);
  \draw[gd8arr] (a.east) -- (i2.west);
  \draw[gd8arr] (a.east) -- (i3.west);
  \node[gd8lab, above=0.2cm of a] {lo stampo};
  \node[gd8lab, above=0.15cm of i1] {le copie nel livello (actor)};
  \node[draw, dashed, fit=(i1)(i2)(i3), inner sep=6pt] {};
\end{tikzpicture}
\caption{Asset e istanze: le modifiche fatte nel livello (nei Details o con il
gizmo) restano sulla singola istanza e non toccano l'asset di partenza.}
\label{gd8:fig:stampo}
\end{figure}

% ---------------------------------------------------------------------
\section{Il transform e la finestra Details}
\label{gd8:sec:transform}

\subsection{Copiare, incollare, annullare}

Prima di parlare del transform, altre scorciatoie ``ereditate'' da Windows:
\begin{itemize}
  \item \textbf{Ctrl+C}, \textbf{Ctrl+V}: copia e incolla. Si crea una copia
        del cubo, che poi si può spostare;
  \item \textbf{Ctrl+clic} su più oggetti: selezione multipla (si vede un
        contorno giallo attorno a tutti gli oggetti selezionati). Con Ctrl+C,
        Ctrl+V si duplicano tutti insieme;
  \item \textbf{Ctrl+Z} ripetuto: si torna indietro fino all'inizio;
  \item per \textbf{rifare} un'azione annullata (\emph{redo}): in alcuni
        programmi è \textbf{Ctrl+Y}, in altri \textbf{Ctrl+Shift+Z}, perché
        nel linguaggio di design di Windows \textbf{Shift inverte
        l'operazione} che si sta facendo. ``È una gran rottura di scatole
        ricordarsi quale ha l'uno e quale l'altro'': se non funziona uno, si
        prova l'altro (Blender, per esempio, usa Ctrl+Shift+Z).
\end{itemize}

\subsection{Il transform}

In termini tecnici, una \textbf{trasformazione} (\emph{transform}) è
l'insieme di tre valori: \textbf{posizione}, \textbf{rotazione} e
\textbf{scala}; in inglese \emph{location}, \emph{rotation}, \emph{scale}.

\subsection{La finestra Details}

Le caratteristiche dell'oggetto selezionato si vedono nella finestra in basso a
destra, \textbf{Details}. Selezionando il cubo, in alto si legge il suo nome
(``Cubo Super Ultra'') e che si tratta di un'\emph{instance}. Le prime tre voci
sono i tre elementi del transform.

\paragraph{Location.} Si può modificare cliccando e trascinando sul campo,
oppure scrivendo il valore: \textbf{X} (rosso), \textbf{Y} (verde),
\textbf{Z} (blu), ``ricordatevelo''. Il docente scrive 400 nella X, preme
Invio, e il cubo si sposta: ``l'ho spostato di 70 centimetri più indietro''
(era a 330).

\paragraph{L'unità di misura.} Nota importante: \textbf{l'unità di misura di
base di Unreal è il centimetro} (si veda anche la
sezione~\ref{ue1:sec:transform}).

\paragraph{Rotation.} Trascinando il primo valore, l'oggetto ruota
\textbf{attorno all'asse X}: l'asse X ``fa da perno''. Il secondo valore lo fa
ruotare attorno all'asse Y, il terzo attorno all'asse Z (quindi
``così'', sul piano orizzontale).

\paragraph{Scale.} È il fattore per cui viene \textbf{moltiplicata} la
dimensione dell'oggetto lungo ciascun asse. Aumentando la scala sulla X,
l'oggetto si allunga lungo la X; lo stesso per Y e Z. Applicato a una sedia
della scena, la si allunga in modo ``psichedelico''. Le tre trasformazioni si
combinano per ottenere l'effetto cercato.

\paragraph{Tutto il resto.} Sotto il transform ci sono tutte le altre
informazioni sull'oggetto selezionato e sulle sue caratteristiche: il
materiale, le caratteristiche fisiche, una casella per \textbf{simulare la
fisica}, una per decidere se l'oggetto \textbf{proietta un'ombra} quando è
colpito da una luce (\emph{Cast Shadow}). Nel corso dell'anno se ne vedrà la
maggior parte, non tutto: alcune voci sono estremamente tecniche e poco
interessanti per il level design. Delle luci ci si occuperà in futuro.

\subsection{W, E, R (e Q)}

Anche il transform ha le sue scorciatoie da tastiera, che cambiano la modalità
del gizmo:
\begin{itemize}
  \item \textbf{W}: spostamento (la modalità di base);
  \item \textbf{E}: rotazione, con un gizmo pensato per ruotare;
  \item \textbf{R}: scala;
  \item \textbf{Q}: solo selezione (``non credo di averlo mai usato, tanto vale
        usare direttamente W'').
\end{itemize}

\paragraph{L'asse appoggiata al muro.} La dimostrazione del docente: per fare
un'asse appoggiata al muro, \textbf{R} per scalare il cubo e renderlo sottile
e lungo, \textbf{W} per avvicinarlo al muro, \textbf{E} per inclinarlo, di
nuovo \textbf{W} per sistemarlo, ``e tac, ho fatto la mia asse appoggiata al
muro''. Poi si tratterà di metterci una texture, e anche questo si vedrà.
``Sembrano tante robe da memorizzare'', ma più si fa pratica, più diventano
automatiche: ``è il bello del cervello umano, siamo in grado di abituarci
molto velocemente''.

\subsection{I gruppi}

Si possono \textbf{raggruppare} più oggetti per spostarli tutti insieme:
selezionarli con \textbf{Ctrl+clic} e premere \textbf{Ctrl+G}. L'esempio:
quando si prepara una libreria e ci si mettono dentro i libri, raggruppare i
libri permette di spostarli, copiarli e incollarli velocemente.

\subsection{Le operazioni nei campi numerici}

Verso fine lezione, mentre la classe lavora, il docente mostra ``un piccolo
trucco magico'' (a voce, perché ha chiuso Unreal per paura che il computer
crolli di nuovo). Tutti i campi numerici dei Details \textbf{accettano
operazioni matematiche} e le eseguono. Se un oggetto ha una scala di 2,75 e la
si vuole triplicare, basta scrivere accanto al valore \texttt{*3} e premere
Invio: il calcolo è automatico. Per dimezzare, \texttt{/2}; per aggiungere 3,
\texttt{+3}. Vale per tutti i campi in cui si può inserire un numero, ed è
così anche in Unity e negli altri engine. ``È molto, molto comodo per
modificare robe con precisione.''

% ---------------------------------------------------------------------
\section{Leggere la documentazione}
\label{gd8:sec:documentazione}

``Nota importantissima'': \textbf{imparate a leggere la documentazione}, e a
interpretare ciò che vi è scritto. ``Quasi tutto quello che vi serve è scritto
da qualche parte.'' Sul sito di Epic Games c'è un'enorme quantità di
informazioni: come funzionano le basi, i contenuti, la creazione di mondi
virtuali, la grafica. Nelle slide c'è un link alla documentazione di Unreal
Engine 5.5, ma si può consultare direttamente quella della 5.8. Il docente
nota, senza commentare oltre, che ``Epic è impazzita per l'AI'': nella
documentazione c'è ``un sacco di roba sull'AI''.

È lo stesso principio della sezione~\ref{gd8:sec:strumento}: saper fare
ricerca e sapere dove cercare è parte dell'imparare a imparare.

% ---------------------------------------------------------------------
\section{I livelli e il salvataggio}
\label{gd8:sec:livelli}

\paragraph{Un livello per ogni zona.} Unreal è diviso in \textbf{livelli}: per
creare una nuova zona, ``tra virgolette'', si crea un nuovo livello, a meno che
non si stia facendo un \emph{open world}. Quello è ``un discorso particolare''
che richiede un lavoro diverso; e anche lì, comunque, ci sono aree chiuse in
cui si entra con caricamenti e simili.

\paragraph{Creare un livello.} \textbf{File} $\rightarrow$ \textbf{New Level},
poi si sceglie un template. Il consiglio è partire da \textbf{Basic}, che ha
già tutto il necessario. \emph{Open World} per ora no; \emph{Empty} nemmeno,
perché per partire da un livello vuoto bisogna aggiungere un po' di cose, e se
ne parlerà più avanti.

\paragraph{Salvare spesso.} ``Abituatevi a salvare spesso.'' In Unreal, un
livello appena creato \textbf{non viene salvato finché non lo si salva
manualmente}. Per salvare:
\begin{itemize}
  \item l'icona del \textbf{floppy disk} (``per chi non sapesse esattamente
        cos'è'': ormai è soprattutto il simbolo del salvataggio);
  \item \textbf{Ctrl+S}: salva il \textbf{contesto attuale} in cui si sta
        lavorando, dato che Unreal può avere più schede aperte in parallelo;
  \item \textbf{Ctrl+Shift+S}: salva \textbf{tutto} (a volte richiede un po'
        di calcolo).
\end{itemize}
In certi casi l'editor chiede dove salvare: si sceglie la cartella del
progetto in cui si vuole mettere il file.

\subsection{Le scorciatoie utili}

Nelle slide c'è un elenco di \textbf{scorciatoie utili}, molte delle quali
corrispondono a funzioni di Unreal che non hanno un pulsante nell'interfaccia e
si possono usare \textbf{solo da tastiera}. Il consiglio è provarle almeno una
alla volta. Quelle viste a lezione sono raccolte nella
tabella~\ref{gd8:tab:scorciatoie}.

\paragraph{Snap to floor.} Una delle più utili è lo \textbf{snap to floor}:
se un oggetto è sospeso in aria e non appoggia sul pavimento, premendo il
tasto \textbf{Fine} (\emph{End}) lo si attacca automaticamente alla superficie
sottostante. (Sul portatile dell'aula il tasto \emph{End} non c'è:
``magnifico''.)

\begin{table}[htbp]
\centering
\small
\begin{tabular}{@{}ll@{}}
\toprule
\textbf{Comando} & \textbf{Effetto} \\
\midrule
Tasto destro + mouse & ruota la telecamera nella viewport \\
Tasto destro + W/A/S/D & sposta la telecamera come in un FPS \\
Tasto destro + Q/E & sposta la telecamera in verticale \\
W / E / R & gizmo di spostamento / rotazione / scala \\
Q & solo selezione \\
F2 & rinomina l'oggetto selezionato \\
Ctrl+clic & selezione multipla \\
Shift+clic (Outliner) & seleziona tutto tra due elementi \\
Ctrl+C / Ctrl+V & copia / incolla (duplica) \\
Ctrl+Z & annulla l'ultima azione \\
Ctrl+Shift+Z (o Ctrl+Y) & ripristina l'azione annullata \\
Ctrl+G & raggruppa gli oggetti selezionati \\
Fine (\emph{End}) & \emph{snap to floor}: appoggia l'oggetto alla superficie sotto \\
Ctrl+S & salva il contesto attuale \\
Ctrl+Shift+S & salva tutto \\
\texttt{*3}, \texttt{/2}, \texttt{+3} in un campo & esegue l'operazione sul valore \\
\bottomrule
\end{tabular}
\caption{Le scorciatoie e i trucchi dell'editor visti nella prima lezione di
level design.}
\label{gd8:tab:scorciatoie}
\end{table}

% ---------------------------------------------------------------------
\section{Il task: il blockout della propria stanza}
\label{gd8:sec:task}

Negli ultimi venticinque minuti circa si parla del lavoro per la volta
successiva, e lo si comincia in classe.

\begin{quote}
\textbf{Creare un blockout della propria stanza}: con le forme di base
(cubi, ma anche cilindri, scale e tutto ciò che c'è nella sezione
\emph{Shapes}), modificate, spostate, allargate, ruotate, ricostruire una
copia della stanza in cui si dorme. Niente colori, niente materiali: ``solo
blocchi''. Consegna: \textbf{tre screenshot} nella cartella dei task, entro
la settimana successiva.
\end{quote}

\paragraph{Perché cominciare in classe.} Il task si comincia subito, non
tanto per finirlo quanto per verificare che tutto funzioni: finestre che non si
vedono, problemi di installazione e simili si risolvono meglio insieme, in
aula. ``Non c'è bisogno che iniziate a fare subito la camera'': si può
cominciare spostando forme e prendendo la mano.

\paragraph{Le domande.}
\begin{itemize}
  \item \emph{Bisogna fare anche pareti, porte, finestre, o bastano
        pavimento e mobili?} ``Come preferite'': è un esercizio molto libero,
        per prendere la mano. L'importante è allenarsi a lavorare
        \textbf{velocemente}: ``metti la cera, togli la cera''. Spostare,
        ridimensionare, ruotare.
  \item \emph{Le forme sono molto geometriche: se volessi smussare i
        cuscini?} Le forme di base \textbf{non si possono modificare}, o
        meglio: si può, ma con uno strumento che si vedrà molto più avanti.
        Si possono però combinare forme diverse. Lo scopo è la pratica con le
        forme e le funzioni di base; chi vuole migliorare il risultato può
        farlo, ma l'importante è dimostrare di aver capito come funziona
        l'engine.
  \item \emph{È il task per la settimana prossima o va consegnato oggi?}
        Per la settimana prossima: lo si è iniziato in classe solo per
        risolvere eventuali problemi tutti insieme.
\end{itemize}

\paragraph{La revisione.} I lavori verranno guardati velocemente la volta
successiva, proiettati sullo schermo dell'aula; ma sarà più una verifica che
il task sia stato fatto. Le correzioni più precise arriveranno con i task
successivi.

\paragraph{Approfondimento: perché la propria stanza.} L'esercizio, oltre a
far prendere la mano con gli strumenti, mette insieme due dei consigli della
lezione: l'\textbf{esperienza sul campo} (osservare lo spazio reale, le sue
misure e le sue proporzioni) e la \textbf{pratica}. Ricostruire uno spazio che
si conosce perfettamente costringe a confrontarsi con la scala: un letto, una
porta, una scrivania hanno misure precise in centimetri, e Unreal ragiona
proprio in centimetri. È anche, in piccolo, il primo blockout: forme di base
che rappresentano oggetti, come il cubo che ``è la casa'' della
sezione~\ref{gd8:sec:cosa-fa}.

% ---------------------------------------------------------------------
\section{Per chi vuole approfondire}
\label{gd8:sec:extra}

Prima di chiudere, ``cinque minuti e trenta secondi di attenzione''. Il level
design è un ambito molto specializzato. Tutti devono imparare almeno ciò che
si fa al primo e al secondo anno: lavorando nei videogiochi si finisce per
parlare con altri designer, e sapere di che cosa si occupano le varie
discipline è importante. Ma è anche innegabile che c'è chi è più interessato
al level design e chi ad altri aspetti, per cui il level design conta meno.
Per i primi, prima della slide del task c'è una \textbf{slide extra} con
materiale aggiuntivo, non obbligatorio.

\paragraph{I talk della GDC.} I primi due link sono video su YouTube:
registrazioni di talk della \textbf{GDC}, la \emph{Game Developers
Conference}\footnote{Il docente la chiama ``Game Developers Convention''; il
nome ufficiale è \emph{Game Developers Conference}.}. È una convention annuale
negli Stati Uniti, considerata ancora oggi \textbf{il punto di ritrovo
dell'industria}. Non è la Milano Games Week, non è il Comic-Con di San Diego:
non ci vanno i giocatori, ci vanno \textbf{gli sviluppatori}, solo chi i
videogiochi li fa. Lì i designer tengono interventi, panel, conferenze e
dibattiti sulla disciplina. Partecipare costa ``veramente tantissimo'' (è la
critica principale che le viene fatta: l'idea è che il biglietto, intorno ai
duemila euro, lo paghi lo studio), ma le registrazioni sono disponibili online.
Si può vedere, per esempio, come gli sviluppatori di \emph{Resident Evil 7}
hanno rinnovato da zero il gioco, e l'intero franchise, raccontato da loro.

\paragraph{Il libro.} C'è poi un link a \textbf{\emph{An Architectural
Approach to Level Design}} (di Christopher W. Totten), ``uno dei migliori libri
sul level design che abbia mai letto'' e un ottimo modo per iniziare ad
affrontare la disciplina. Lo si trova a molto meno dei 67 dollari indicati nel
link; il docente lo ``stra-consiglia''.

Non sono obbligatori: chi è incuriosito li trova lì, soprattutto i talk; chi
preferisce specializzarsi in altro non ha problemi. Alcuni degli argomenti
trattati in quei talk e in quel libro verranno comunque affrontati nel corso
dell'anno.

% ---------------------------------------------------------------------
\section{La chiusura}
\label{gd8:sec:chiusura}

Come sempre, qualche minuto per domande ``varie ed eventuali'', anche se alla
prima lezione di solito non c'è molto da chiedere; e per una volta, nota il
docente, Unreal ha funzionato a tutti (tranne a lui). Uno studente ha già
quasi finito la stanza, e chiede chiarimenti sulle forme geometriche (si veda
la sezione~\ref{gd8:sec:task}).

\paragraph{Come contattare il docente.} Lo si trova su \textbf{Discord},
cercando ``Stefano Zocchi'' nel server della scuola\footnote{La trascrizione
riporta il nome del server come ``Hive''; non è verificabile.}, oppure per
email, che il docente comunica alla classe per le questioni più complesse.

\paragraph{Il riepilogo.} ``Per la prossima volta voglio che mi creiate una
simulazione di camera vostra fatta solo con le forme di base, senza colori,
niente, solo blocchi.'' Non devono essere per forza cubi: cilindri e tutte le
altre forme della sezione \emph{Shapes}. ``Divertitevi, usate la fantasia.''
Tre screenshot nella cartella dei task. Ci si rivede la settimana successiva.

% ---------------------------------------------------------------------
\section*{Punti chiave del capitolo}
\addcontentsline{toc}{section}{Punti chiave del capitolo}

\begin{itemize}
  \item Il corso di \textbf{Level Design} unisce teoria, metodologia e
        soprattutto \textbf{pratica in Unreal Engine 5}, con pratica assistita,
        task a casa (da rivedere in classe) ed esercizi teorici.
  \item Un livello brutto va analizzato: \textbf{perché} non funziona e
        \textbf{come} lo si migliorerebbe. La Cattedrale del Sangue Sacro di
        \emph{Code Vein} riassume gli errori tipici: labirinto, piani non
        collegati, rumore visivo, nessun landmark, ripetitività, nessun
        interesse nell'esplorare. Bisogna anche distinguere
        \textbf{scelte creative} (la Città Infame, \emph{Nightreign},
        \emph{Frostpunk 2}) da \textbf{errori}, e valutare le conseguenze di
        ogni modifica.
  \item \textbf{Parlare in pubblico} è indispensabile: un designer mediocre
        che sa presentare le proprie idee avrà più successo di un designer
        geniale che non ci riesce.
  \item Il \textbf{level design} è la disciplina che crea livelli, ambienti,
        spazi di gioco, di solito con un \textbf{level editor}; ha un lato
        artistico e uno tecnico, e altri nomi (environment design, game
        mapping). Il \textbf{blockout} usa forme di base prima dell'arte
        finale.
  \item Il level designer \textbf{crea spazi}, ma ``level design'' è un
        \textbf{termine ombrello}: blockout, luci (il \emph{lighting pass},
        come nel cinema), set dressing, documentazione (layout, mappe,
        diagrammi di flusso), test iterativo (playtest, QA), narrazione
        ambientale (dove piazzare l'informazione cambia l'esperienza).
  \item La parola magica è \textbf{``dipende''}. L'engine dipende
        dall'azienda (Unreal, Unity, Godot, GameMaker, engine proprietari):
        lo strumento vero del level designer è \textbf{il level designer},
        e la capacità più importante è \textbf{imparare a imparare}.
  \item Per imparare: \textbf{teoria} (vale ovunque; il \textbf{landmark} è
        un elemento distinguibile e visibile da lontano che aiuta a
        orientarsi; anche rifiutarlo dev'essere una scelta consapevole),
        \textbf{pratica}, \textbf{expertise}, \textbf{esperienza sul campo},
        \textbf{lavoro di gruppo}. E soprattutto pratica.
  \item Il \textbf{picco della fiducia} (tra primo e secondo anno) e
        l'\textbf{effetto Dunning-Kruger} (meno si sa, più ci si sente
        sicuri): l'obiettivo è attraversare la valle senza abbandonare.
  \item Tre corollari: \textbf{non esistono regole fisse} (isolare la
        domanda e la ragione della scelta); il design è \textbf{olistico}
        (conoscere almeno le basi di tutte le discipline); il design è
        \textbf{artigianato}, non arte (ha uno scopo). Fare giochi, ricreare,
        trovare una comunità, farsi ispirare.
  \item Editor: \textbf{viewport} (tasto destro + WASD, Q/E), \textbf{Place
        Actors} (Shapes), \textbf{gizmo} (frecce, quadratini, centro; il
        pivot lo decide chi fa la mesh), \textbf{Outliner} (ordine
        alfabetico, cartelle, F2), \textbf{Content Browser}, \textbf{Details}.
  \item \textbf{Asset} = file nel progetto (lo stampo); \textbf{actor} =
        istanza nel livello (la copia). Le modifiche fatte nel livello restano
        sull'istanza.
  \item \textbf{Transform} = location, rotation, scale; l'unità è il
        \textbf{centimetro}. W/E/R per spostare, ruotare, scalare; Ctrl+G per
        raggruppare; \emph{End} per lo snap to floor; i campi numerici
        accettano operazioni (\texttt{*3}, \texttt{/2}).
  \item \textbf{Leggere la documentazione}. Nuovo livello: File, New Level,
        Basic; \textbf{salvare spesso} (Ctrl+S, Ctrl+Shift+S).
  \item Task: il \textbf{blockout della propria stanza} con le sole forme di
        base, tre screenshot nella cartella dei task. Per approfondire: i talk
        della \textbf{GDC} e \emph{An Architectural Approach to Level Design}.
\end{itemize}
